\documentclass[11pt]{article}
\usepackage[a4paper,margin=2cm]{geometry}
\usepackage{amsmath,amssymb,amsfonts}
\usepackage[T1]{fontenc}
\usepackage[utf8]{inputenc}
\usepackage{lmodern}
\usepackage{microtype}
\usepackage{booktabs}
\usepackage{tabularx}
\usepackage{enumitem}
\usepackage[hidelinks]{hyperref}
\usepackage{xcolor}

\DeclareMathOperator*{\argmin}{arg\,min}
\DeclareMathOperator*{\argmax}{arg\,max}

% --- density settings -----------------------------------------------------
\setlength{\parindent}{0pt}
\setlength{\parskip}{4pt plus 1pt minus 1pt}
\setlength{\abovedisplayskip}{5pt plus 2pt minus 2pt}
\setlength{\belowdisplayskip}{5pt plus 2pt minus 2pt}
\setlength{\abovedisplayshortskip}{2pt plus 1pt}
\setlength{\belowdisplayshortskip}{2pt plus 1pt}
\setlist{nosep,leftmargin=2em}
\allowdisplaybreaks

% --- helper environments for plain-language notes and examples -------------
\newenvironment{plainword}{%
  \par\medskip\noindent\begin{minipage}{\textwidth}%
  \color{gray!75!black}\small\itshape
}{%
  \end{minipage}\par\medskip}

\newenvironment{example}{%
  \par\smallskip\noindent\begin{minipage}{\textwidth}%
  \color{teal!60!black}\small
  \textbf{\color{teal!70!black}Real-world example.}\
}{%
  \end{minipage}\par\smallskip}

\begin{document}
\begin{center}
{\LARGE\bfseries Relational Latent World Learning\par}
\end{center}

\begin{center}
{\large An Interaction-Centred Research Position on Effective Variables, Task Constraints, and Action\par}
\end{center}

\vspace{0.6em}
\begin{plainword}
This paper describes a research plan, not finished results. Its one big question is:
\emph{can a machine (or an agent) figure out how a part of the world works simply by trying things and watching what happens},
and then \emph{reuse} that understanding both to keep learning and to actually get things done? The words are kept
simple below, and every abstract idea is followed by a concrete example --- mostly from medicine, because
"observing a system, acting on it, and learning from the response" is exactly what a doctor does every day.
\end{plainword}

\subsection*{How to read this paper (a short glossary)}

Throughout the paper we use a small set of technical words. Here they are in plain language, each with a medical example.

\begin{center}
\small
\begin{tabularx}{\textwidth}{@{}p{0.24\textwidth}p{0.38\textwidth}X@{}}
\toprule
Term & Plain meaning & Medical example \\
\midrule
Agent & Something that can see and act in a situation & A robot arm, or a doctor \\
Observation $O_t$ & What you can see or measure at one moment & Vital signs, a blood test, an X-ray \\
Action $a_t$ & What you actually do & Giving a drug, pushing an object \\
State $S_t$ & A short useful summary of "where things stand now" & "Blood pressure high, kidneys failing" \\
Transformation (TP) $p$ & \emph{What actually changed}, not the action that caused it & "Blood pressure fell by 20\,mmHg" (not "I gave lisinopril") \\
Mechanism $Z_D$ & A reusable cause--effect unit that recurs & "This drug lowers BP via the kidney" \\
Relation $R_D$ & A one-way effect of one mechanism on another & "Drug A makes drug B stronger" \\
World $W_D$ & The whole learned "how this domain works" model & A doctor's understanding of cardiology \\
Task $C_T$ & A goal, a desired outcome & "Get this patient's BP to normal" \\
Exploration & Choosing the next action to reduce confusion & Ordering a test to tell two diagnoses apart \\
Identifiability & Whether the data can tell two explanations apart & "Is it A or B? The data must decide." \\
\bottomrule
\end{tabularx}
\end{center}

\begin{plainword}
The key move of the whole paper is this: the agent is allowed to describe what happened in terms of \emph{reusable mechanisms}
and \emph{relations between them}, and is only allowed to keep a mechanism or a relation if it keeps being useful on new situations.
A description is not accepted just because it looks pretty --- it must \emph{earn its place} by predicting things that were never shown to it.
\end{plainword}

\subsection*{Abstract}

This paper asks whether an agent that can sense and act can discover a compact, relational account of a domain from
interaction alone, and then reuse that account for two jobs at once: exploring on its own, and completing tasks.
The proposed answer is neither a complete hand-written physics engine nor a black-box task robot. It is a stack of
\emph{useful summaries}, each sitting on top of the last.

In plain terms, the stack is built like this:

\begin{itemize}
\item A \emph{persistent state} $S_t$ makes "change" measurable, so we can compare one moment to the next.
\item A \emph{transformation primitive} (TP) $p_i$ describes \emph{what happened} over a short stretch of time --- the change, not the action.
\item \emph{Reusable domain variables} $Z_D$ and \emph{directed relation operators} $R_D$ explain the recurring parts of those changes.
\item Together these form a \emph{relational latent world} $W_D=(Z_D,R_D)$: a compact model of "how this domain behaves."
\end{itemize}

\begin{example}
A doctor does not memorize one unrelated fact per patient. She builds a reusable model: "this class of drugs lowers blood
pressure through the kidney," "this second drug makes the first stronger." Those are mechanisms and relations. When a new
patient arrives, she reuses the same model instead of starting from zero.
\end{example}

This world closes a learning loop: when it is confident, it asks for missing evidence; when it is not, it asks to try
transformations at the edge of what it currently understands. In parallel, several demonstrations of the same task are used to
learn a \emph{task constraint} $C_T$: a compact description of the goal, learned by predicting one demonstration from another.
The task state is deliberately minimal, $Z_T=(S_t,C_T)$. The task is not stored as a new object; instead $C_T$ \emph{acts on} the
explicit world $W_D$ to produce a distribution over "next transformations that fit the task," plus a "should I stop now" signal.
Low-level control is then obtained by asking the already-learned model the reverse question: \emph{"which action would produce
that transformation?"}

There is no separate desired-transformation variable $D_T$, no explicit task graph $G_T$ in the first implementation, no learned
uncertainty object $U_D$, and no separate exploration or task motor policy. The paper keeps these out deliberately, and says why.

The paper also places this proposal against related ideas --- object-centric learning, equivariant representations, temporal
abstraction, causal representation learning, relational inductive bias, active experiment design, task inference, imitation
learning, and model-based control. For each, it states plainly what is borrowed, what is changed, and why existing guarantees do
not carry over. The final product is a \emph{falsifiable} research programme: world variables must earn their status through
reusable predictions; task context must earn its status through transfer across demonstrations; and every abstraction is judged
by its held-out consequences, not by whether its learned "shape" looks clean.

\section*{1. Scientific Problem Definition}

\subsection*{1.1 The scientific question}

Let an agent observe $O_t$, carry out a low-level action $a_t$, and receive the next observation $O_{t+1}$. The scientific question is:

\[
\boxed{
\begin{minipage}{0.88\textwidth}
Can the agent use time-aligned observation--action streams to discover reusable effective variables and their directed
relations, actively gather the evidence needed to refine that structure, infer a persistent task constraint from different
demonstrations of the same task, and turn that constraint back into actions --- \emph{without ever being told a human
vocabulary} of objects, mechanisms, or named actions?
\end{minipage}}
\]

\begin{example}
The agent is never told "this is a drug, this is a kidney, this is 'blood pressure.'" It only sees: before (the numbers),
what was done (the dose), and after (the new numbers). From many such triples it must discover that there is a stable thing
called "blood pressure," that a certain intervention reliably changes it, and that this change can be undone or boosted by
other interventions.
\end{example}

The primitive unit of data is a continuous interaction trajectory

\[
\boxed{\tau=(O_0,a_0,O_1,a_1,\ldots,a_{T-1},O_T),}
\]

with strict alignment in time between what was observed and what was actually done.

The proposal rejects two shortcuts. First, an \emph{action label} is not the same as the \emph{transformation} it produces: the
same action can have different effects in different situations, and different actions can produce the same effect. Second, a
\emph{task label} or a \emph{final picture} is not a sufficient description of a task: different strategies can satisfy the same
goal, while identical endpoints can hide different required processes.

\begin{example}
"Push the cup" is one action, but the \emph{change} depends on the cup: an empty cup slides, a full cup tips, a glued cup does
not move. Conversely, "slide the cup left" can be achieved by pushing from the right or pulling from the left. So what the agent
should record is the \emph{change}, not the name of the action. Likewise, "blood pressure returned to normal" can be reached by
a diuretic, by a beta-blocker, or by fluid restriction --- the endpoint alone does not tell you the process.
\end{example}

\subsection*{1.2 Target explanatory hierarchy}

The proposed hierarchy is

\[
\boxed{
O_t\rightarrow S_t\rightarrow TP\rightarrow
\left\{
\begin{array}{l}
Z_D\leftrightarrow R_D\rightarrow W_D\rightarrow \text{exploration},\\
TP\text{ sequences}\rightarrow C_T\rightarrow \text{task transformation}
\end{array}
\right.
\rightarrow \text{intervention}.}
\]

\begin{plainword}
Reading left to right: raw observation $\to$ a persistent state $\to$ a "what changed" summary (TP). That summary then feeds two
branches: (1) reusable mechanisms and their relations form a world model, used for exploration; (2) sequences of changes are used
to learn a task goal. Both branches end back in a concrete action. The whole loop is "see, act, learn, act better."
\end{plainword}

The terms play deliberately different roles:

\begin{itemize}
\item $S_t$ is a persistent coordinate system for interaction-relevant state.
\item $p_i$ is an experience-level description of one concrete local change over a variable-length stretch of time.
\item $Z_D=\{Z_D^j\}$ is a domain-level bank of reusable effective mechanisms; $z_i^j$ is the value of mechanism $j$ in experience $i$.
\item $R_D=\{R_D^{j\rightarrow k}\}$ is a set of stable directed operators describing how one mechanism changes another.
\item $W_D=(Z_D,R_D)$ is the shared relational latent world for the chosen domain and scale.
\item $C_T$ is a persistent task constraint inferred from a TP sequence and trained to transfer across different demonstrations of the same task.
\item $Z_T=(S_t,C_T)$ is the minimal task state. $W_D$ stays an external, domain-shared structure rather than being folded into $Z_T$.
\end{itemize}

Task transformation is \emph{not} a new latent variable. It is the operation

\[
\boxed{
\mathcal T_T:(S_t,W_D,C_T)\longmapsto
P_T(p_{\mathrm{next}},y^{\mathrm{stop}}\mid S_t,W_D,C_T),}
\]

possibly also looking at the most recent TP as a short-term "progress cue." In plain terms, $C_T$ does not create a new world; it
\emph{re-weights} the transformations that $W_D$ already makes possible, keeping the ones that fit the goal. The first
implementation therefore has neither a stored "desired next transformation" $D_T$ nor a separately built task graph $G_T$.

\subsection*{1.3 Why this problem is not ordinary representation learning}

Predicting the next observation is not enough to pin down the right concepts: many different internal coordinate systems can
reproduce the same one-step predictions. So this paper asks a stronger, functional question. Which variables keep explaining
effects \emph{across} contexts? Which directed relations stay necessary under controlled comparisons? Which task code changes
future transformation predictions when the world state is held the same? A latent variable is accepted not because it looks
interpretable or statistically separated, but because it \emph{earns operational meaning} by being reusable.

\begin{example}
A student could memorise "this exact patient, this exact pill, this exact result" --- that predicts each case perfectly but
generalises to nothing. What we want instead is the reusable rule "ACE inhibitors lower blood pressure," which transfers to new
patients, new doses, and new contexts. The test of a concept is \emph{does it keep working somewhere it has never been?}
\end{example}

This also changes the role of action. Action is not merely a control target. It is \emph{experimental leverage}. Different actions
supply the contrasts that passive watching cannot, and the learned structure in turn tells the agent which action would be
informative next. The world model is therefore both an explanation and an experiment-design state.

\subsection*{1.4 Scope and non-claims}

The framework targets a bounded domain, sensor suite, action space, spatial scale, temporal scale, and abstraction level. It does
not claim to recover fundamental physical variables, a unique object ontology, or a universally valid causal graph. $Z_D$ is an
\emph{effective} description relative to what the agent can observe and change. Relations are functional modulation operators in
latent effect space, not automatically causal effects in the formal interventionist sense.

Four kinds of statements are kept strictly separate throughout:

\begin{enumerate}
\item \textbf{Physical or structural prior:} a claim about regularity in the domain or task family.
\item \textbf{Identifiability condition:} a condition on the data needed to tell hypotheses apart.
\item \textbf{First implementation:} an architectural, distributional, or optimisation choice.
\item \textbf{Diagnostic:} an evaluation or evidence quantity that is not promoted into a latent object or a strong training loss.
\end{enumerate}

This distinction matters. A Transformer, a Gaussian posterior, a Hard-Concrete relaxation, pairwise message passing, task
grouping supplied by a dataset, or local differentiability of a learned controller is \emph{not} a physical assumption.

\begin{example}
"We used a Transformer" is a statement about our implementation, like "we wrote the notes in this notebook." It says nothing
about whether the body actually works the way our model imagines. Keeping these two kinds of claim apart is what stops a good
prediction from being mistaken for a true explanation.
\end{example}

\section*{2. Literature Review and Research Positioning}

\begin{plainword}
This section says, for each neighbouring field: "here is the useful idea we take, here is the important way we differ, and here
is a guarantee from that field that does \emph{not} carry over to us." You can read just the first sentence of each subsection to
get the idea; the detail is there if you want it.
\end{plainword}

\subsection*{2.1 Object-centric and temporally persistent representations}

Slot Attention introduced exchangeable "slots" that bind perceptual features into a set of task-dependent representations [1].
SAVi added temporal conditioning for video, showing that earlier slots can guide binding in later frames [2]. SOLD showed the value
of slot-structured latent dynamics for relational manipulation and model-based reinforcement learning [3]. We borrow the
set-valued latent interface, temporal conditioning, and action-conditioned dynamics. We differ in ontology and purpose: a slot
here is a persistent \emph{interaction unit}, not a claim that the world is made of objects, and slots are only the coordinate
system from which transformations and mechanisms are discovered. Consequently, object-discovery results, segmentation quality, or
empirical control gains do not guarantee TP, $Z_D$, or $R_D$ identifiability. Slot permutation, splitting, merging, and non-object
solutions remain possible.

Vision Transformers provide a convenient localised token interface [4], and Set Transformers and Deep Sets provide
permutation-aware computation over sets [18,19]. These architectures are borrowed as implementation tools. No theorem about the
world follows from choosing them, and no representation guarantee transfers from their universal-approximation or attention
properties.

\subsection*{2.2 Invariance, equivariance, and interaction-grounded change}

Equivariant representation learning formalises "predictable latent change under known transformation groups." SIE explicitly
splits invariant and equivariant components and stops equivariant information from collapsing into the invariant branch [29].
Symmetry-based approaches clarify an important goal: irrelevant variation should be thrown away while action-relevant change stays
predictable. We borrow this separation of stable content from structured change. We differ because robot interventions need not
form a known group, be invertible, or act globally; contact, containment, irreversible change, and state-dependent effects all
violate the clean symmetry setting. The proposed TP is therefore learned from \emph{realised} state change and executed actions,
not from supervised relative group elements. Group-equivariance, losslessness, and disentanglement guarantees do not transfer.

This difference is substantive. The target is not "invariance to action" but a representation in which

\[
\boxed{(S,A)\mapsto p\mapsto \Delta S}
\]

is stable enough to support mechanism discovery and inversion. What stays invariant is the reusable \emph{mechanism identity};
what varies (equivariantly, in a loose operational sense) is its \emph{realisation and effect}.

\begin{example}
The mechanism "this drug dilates blood vessels" is the same thing across patients (the invariant identity). But \emph{how much}
dilation happens varies with dose, weight, and other drugs present (the varying realisation). We want a representation that keeps
the first fixed and lets the second vary.
\end{example}

\subsection*{2.3 Temporal abstraction, skills, and transformation primitives}

The options framework formalised temporally extended actions in semi-Markov decision processes [5]. Temporal variational
inference and later temporally abstract world models learn latent skills or options from trajectories [6,30]. We borrow variable
duration, learned termination, and the principle that temporal segmentation should reflect a coherent latent regime. We differ in
what the segment's latent \emph{denotes}. A skill or option is normally a policy-side object --- \emph{how to act} --- whereas a TP
denotes the realised world-side transformation --- \emph{what happened}. The TP posterior therefore excludes action, while its
prospective prior includes state and intervention. Option-value, policy-composition, and semi-MDP guarantees do not establish that
the inferred TP is action-invariant, mechanism-reusable, or causally meaningful.

\begin{example}
A "skill" would be "the manoeuvre for inserting an IV" (how to act). A "transformation" would be "the needle is now inside the
vein" (what changed). The same transformation can be reached by different skills, and the same skill can fail to produce the
transformation in a different body. We model the transformation, because that is what both learning and task success depend on.
\end{example}

Latent dynamics and predictive world models show that useful representations can be learned through future prediction and used for
planning [3,26--28,31]. Predictive invariant/equivariant world models further motivate separating stable content from predictable
change [7]. Latent-action work shows both the promise of discovering action-like coordinates and the risk that distractors or
missing intervention information contaminate them [8,9]. We borrow action-conditioned latent prediction and receding-horizon
optimisation, but keep the executed intervention explicit and distinguish the \emph{effect} posterior from the
\emph{intervention-conditioned} prior. We also insert explicit transformation, mechanism, and relation levels, and attribute errors
across them. Predictive accuracy or latent-action cycle consistency alone does not identify the proposed hierarchy; a monolithic
model can fit transitions while hiding exactly the structure we care about.

\subsection*{2.4 Disentanglement and causal representation learning}

Unsupervised disentanglement is non-identifiable without inductive biases on the model and data [32]. Nonlinear ICA becomes
identifiable only with additional structure such as observed auxiliary variables [33]. These negative and positive results motivate
the explicit statement of priors and coverage conditions in this paper. We do not claim that sparsity or prediction uniquely
recovers true factors.

CITRIS uses temporal sequences with observed intervention targets to identify multidimensional causal factors under its generative
assumptions [10]. BISCUIT studies unknown binary interaction indicators that switch causal mechanisms [11]. Mechanism-sparsity work
shows that sparse latent dependencies or intervention masks can yield identifiability under graph and data conditions [12].
Multi-node and linear-system intervention analyses relax or clarify intervention requirements in more restricted mixing settings
[13,14], while interventional causal representation learning derives identifiability from support changes under interventions [15].
iCITRIS extends temporal causal representation learning to instantaneous effects [16]. We borrow four principles: temporal
information, intervention diversity, sparse or structured mechanism participation, and the need to state what variation identifies
an edge.

The differences are equally important. The present model uses learned slots, variable-duration TPs, unknown many-to-many mechanism
allocation, nonlinear mechanism modules, and directed residual relation operators. Intervention targets are not supplied. $Z_D$ is
effective and domain-relative, not assumed to equal a ground-truth causal variable. The learned relation graph need not be acyclic.
Hence the identifiability theorems of [10--16,33] do not transfer. They motivate diagnostics and experimental designs, not a proof
that this neural hierarchy recovers a unique causal model.

\begin{example}
Most causal-representation results assume someone has \emph{labelled} which variable was intervened on ("I changed the dose of
this one drug"). Here the agent gets no such labels --- it only sees that \emph{something} was done and \emph{something} changed. So
the existing guarantees, which assume labels, cannot simply be imported.
\end{example}

\subsection*{2.5 Reusable mechanisms, latent feature allocation, and relational structure}

The Indian Buffet Process supplies a conceptual precedent for sparse, many-to-many latent feature allocation [17]. We borrow the
idea that one experience may activate several reusable features and one feature may recur across many experiences. We do not use the
IBP generative process or its nonparametric inference guarantees; the first implementation uses an overcomplete finite bank and
differentiable participation gates.

Graph networks and relational inductive biases motivate representing entities and their interactions explicitly [22]. FiLM supplies
a simple conditioning operator [21], and $L_0$ regularisation supplies differentiable sparse gates [20]. We borrow directed message
passing, constrained modulation, and sparse structural selection. We differ in edge semantics: nodes are reusable effective
mechanisms and an edge is the residual modulation of a target mechanism by a source mechanism. The graph is not an input object
graph, and graph-network expressivity does not make an edge identifiable. Direction requires matched-context contrasts and held-out
explanatory power.

\subsection*{2.6 Active learning, experimental design, and curiosity}

Bayesian experimental design defines experiment value through expected information gain [23], while D-optimal and log-determinant
criteria reward measurements that expand rather than duplicate supported directions [24]. Curiosity methods reward prediction error
as an intrinsic signal [25]. We borrow the principle that data collection should be driven by expected reduction of structural
ambiguity, and use an explicit log-determinant coverage statistic as a first estimator.

We differ from generic curiosity in two ways. First, novelty is not intrinsically valuable: evidence must bear on mechanism reuse
or directional relation identification. Second, disagreement is rewarded only when the lower TP level is locally competent;
otherwise the agent would chase sensor noise or exploit its own predictor. We also do not maintain a full Bayesian posterior over a
neural world model. Therefore Bayesian optimality, submodularity, and regret guarantees do not transfer to the learned nonlinear
objective.

\begin{example}
A curious student who is rewarded for "anything surprising" will waste time being surprised by noise. A good doctor is surprised
only by things that change the diagnosis or the plan. We want the second kind of curiosity: surprise that is \emph{useful}.
\end{example}

\subsection*{2.7 Task representations from demonstrations}

One-shot visual imitation and Task-Embedded Control Networks infer a task representation from demonstrations and condition a
controller on it [34,35]. Probabilistic-context meta-RL separates task inference from control [36]. Conditional Neural Processes
separate a context set from target prediction [37]. We borrow the context/query logic: infer $C_T$ from one demonstration and
require it to predict \emph{another} demonstration of the same task. This is stronger than reconstructing its own trajectory,
because it discourages a path-specific code.

The proposed task variable differs from prior task embeddings in both supervision and interface. The first experiment receives
episode boundaries and same-task grouping, but \emph{not} semantic task labels as inputs to the network. $C_T$ predicts
transformations rather than direct motor commands or rewards. It acts on the explicit $W_D$, and the low-level action is found by
TP-prior inversion. Meta-learning generalisation, imitation success, or neural-process consistency does not imply that $C_T$ has
captured a \emph{task constraint} instead of strategy, initial state, demonstrator identity, or dataset bias. That claim must be
tested with cross-strategy demonstrations, cross-initial-state transfer, and a no-$C_T$ predictive ablation.

\begin{example}
If two doctors both "reduce blood pressure" but one uses a diuretic and the other uses diet, a good task code should capture the
shared \emph{goal} ("lower BP"), not "which doctor" or "which drug." We test this by asking the code learned from one doctor's
demonstration to predict the other doctor's demonstration.
\end{example}

\subsection*{2.8 Multimodal task progression and model-based action realisation}

Demonstrations of the same task can contain several valid next transformations. A unimodal regressor can average incompatible modes
and pressure the task encoder to leak strategy identity. Mixture-density models provide a direct probabilistic representation of
multimodal conditional outcomes [38]; diffusion policies show the practical importance of multimodal action distributions in robot
imitation [39]. We borrow only the need for a multimodal output and choose a conditional Gaussian mixture in TP space for the first
implementation. Diffusion-policy control performance and density-model expressivity do not guarantee that the modes correspond to
meaningful task alternatives.

\begin{example}
"Lower blood pressure" can next be achieved by "more diuretic" or by "add a beta-blocker." Averaging the two would produce a
meaningless half-measure. A mixture model keeps both options alive as separate choices, which is what real medicine (and real
robotics) requires.
\end{example}

Differentiable MPC, PETS, visual foresight, and Universal Planning Networks optimise actions through learned predictive models
[26--28,40]. We borrow direct action-sequence optimisation, uncertainty-aware or sampling baselines, and receding-horizon
replanning. The difference is the objective: exploration maximises structural evidence and abstraction challenge; task execution
matches the TP distribution demanded by $\mathcal T_T$. No independent motor policy is learned. Classical MPC guarantees do not
transfer because the learned TP and world abstractions are approximate, stochastic, and only locally calibrated.

\subsection*{2.9 Positioning summary}

\begin{center}
\small
\begin{tabularx}{\textwidth}{@{}p{0.19\textwidth}XXX@{}}
\toprule
Line of work & Borrowed or inspired & Distinctive move here & Guarantee that does not transfer \\
\midrule
Object-centric video & persistent set-valued state & slot as interaction unit, not object ontology & object recovery or slot identifiability \\
Equivariance & separate stable identity from structured change & interventions need not be known groups & group-equivariance and losslessness \\
Temporal abstraction & variable duration and termination & TP represents realised effect, not a policy option & semi-MDP or option-control guarantees \\
Causal representation learning & interventions, sparsity, contrast, identifiability discipline & unknown targets, learned TP mediator, effective mechanisms, possibly cyclic relations & recovery of true causal variables/graph \\
Graph networks & explicit directed interactions & edge is source-to-target mechanism modulation & edge semantics or direction identifiability \\
Experimental design & non-redundant evidence gain & structure-specific evidence plus TP-gated abstraction frontier & Bayesian optimality or submodularity \\
Task/meta learning & demonstration-derived context & cross-demonstration TP prediction on explicit $W_D$ & strategy-invariant task identity \\
Model-based control & optimise actions through predictions & invert TP prior for both exploration and tasks & stability or optimality under model error \\
\bottomrule
\end{tabularx}
\end{center}

\section*{3. Physical and Structural Priors}

\begin{plainword}
A "prior" here means an assumption we make \emph{before} looking at data --- a belief about how the world is organised. We list only
assumptions about the domain itself here. Choices about networks, distributions, and optimisers come later. Each prior gets a
plain-language restatement and an example.
\end{plainword}

\subsection*{P1. Domain-local effective compressibility}

At a specified domain, spatial scale, temporal scale, and abstraction level, intervention-relevant change can be explained by fewer
reusable effective mechanisms than raw observation dimensions:

\[
\boxed{\Delta S\approx F_W(S;Z_D,R_D),\qquad |Z_D|\ll \dim(O).}
\]

\begin{plainword}
A small number of reusable causes can explain most of what changes, even though the raw measurements (pixels, readings) are huge in
number.
\end{plainword}

\begin{example}
A patient's full monitoring trace has thousands of numbers, but a doctor summarises "the blood pressure is dropping because of
these two causes." The summary is short, but it explains the change. That is what P1 assumes: the domain \emph{allows} such short
summaries.
\end{example}

The mechanisms are not asserted to be fundamental physical variables or unique across domains.

\subsection*{P2. Intervention--response regularity}

The effect of a low-level intervention depends systematically on the pre-intervention state:

\[
\boxed{P(S_{t+1}\mid S_t,a_t)\neq P(S_{t+1}\mid a_t)\ \text{in general}.}
\]

\begin{plainword}
What an action does depends on where you are. The same action in different states generally gives different outcomes.
\end{plainword}

\begin{example}
Giving a diuretic to a dehydrated patient is very different from giving it to an overloaded one. The drug is the same; the state is
not. So "action $\to$ outcome" is not a clean fixed rule --- the state matters.
\end{example}

Comparable transformations can recur across different controls, and comparable controls can yield different transformations under
different states. This regularity is what makes a transformation variable distinct from an action code.

\subsection*{P3. Temporally local transformation regimes}

Over some bounded intervals, a single latent transformation description is sufficient to explain the local evolution better than
re-describing the code at every timestep. Boundaries need not coincide with action switches and need not be perfectly discrete in
every domain. The prior is local persistence, not a preferred duration.

\begin{example}
"Steadily lowering the blood pressure over half an hour" is one coherent transformation, even though it is implemented by many
small titration steps. It would be wasteful to describe each small step as a brand-new phenomenon.
\end{example}

\subsection*{P4. Reusable and locally compositional mechanisms}

The same mechanism function can recur across contexts and tasks, while a local TP normally engages a subset of the mechanisms
available in the domain:

\[
\boxed{e_i^j=\mathcal M_j(S_i^{\mathrm{pre}},v_i^j),\qquad
\mathbb E\|m_i\|_0<M_{\max}.}
\]

\begin{plainword}
Each mechanism is a reusable function. In any single change, only a \emph{few} mechanisms are usually active, and one change can
activate several at once (many-to-many).
\end{plainword}

\begin{example}
The mechanism "renin--angiotensin--aldosterone axis" is reused across many patients. A given clinical event might activate this
mechanism together with "sympathetic activation" --- two mechanisms, one event. And the same mechanism shows up in thousands of
different events. That reuse is what P4 assumes.
\end{example}

This does not impose one TP per mechanism. The mapping is many-to-many.

\subsection*{P5. Stable directed relational modulation}

Some effects cannot be explained by isolated mechanisms alone but admit reusable, directed source-to-target modulation:

\[
\boxed{R_D^{j\rightarrow k}\neq R_D^{k\rightarrow j}\ \text{in general}.}
\]

\begin{plainword}
Some effects need two mechanisms plus a \emph{direction} between them. "A changes B" is not the same as "B changes A."
\end{plainword}

\begin{example}
Drug A may increase the effect of drug B (A$\to$B), while B does nothing back to A. The arrow has a direction. Whether the current
data actually \emph{support} that arrow is a separate question --- P5 only assumes such directed relations can exist.
\end{example}

Relation existence is a property of the domain-level mechanism system; whether the current dataset supports it is an epistemic
question.

\subsection*{P6. A shared domain world under task variation}

Within the chosen domain, tasks normally change \emph{which} feasible transformations are relevant, not the underlying domain
mechanisms themselves. Thus $W_D$ is shared, while $C_T$ reweights its currently realisable transformation possibilities. This prior
motivates freezing the world representation when first learning task context.

\begin{example}
Cardiology (the domain) does not change between the goal "lower blood pressure" and the goal "raise blood pressure." The same
mechanisms apply; the goals just pick different targets. So we learn the world once, and let each task re-weight it.
\end{example}

\subsection*{P7. Cross-realisation persistence of task constraint}

Different successful trajectories can instantiate the same task constraint. Therefore a task has information that persists across
variation in initial state, strategy, and low-level execution:

\[
\boxed{\tau_{g,r}\sim g,\ \tau_{g,s}\sim g
\quad\Rightarrow\quad
C_T^{g,r}\ \text{and}\ C_T^{g,s}\ \text{are functionally interchangeable}.}
\]

\begin{plainword}
Two different successful attempts at the same task should give task codes that do the same job, even if their numbers differ.
\end{plainword}

\begin{example}
Two different patients, two different doctors, two different drug regimens --- but the same goal "restore normal blood pressure."
The task code learned from one should be able to predict sensible progress in the other. "Functionally interchangeable" means
each code predicts the other's progression; it does not require identical coordinates.
\end{example}

\subsection*{P8. Task relevance is expressible in transformation space}

For the first target domains, the information needed to tell task-consistent from task-inconsistent progress can be expressed as a
distribution over TPs and termination conditional on the current world state. This is a genuine scope prior. If two tasks demand
different information that never changes their TP progression or stop condition, the proposed $C_T$ interface cannot distinguish
them.

\subsection*{What is deliberately not a physical or structural prior}

The following are implementation or experimental choices: ViT features, slots, Transformers, FiLM, diagonal Gaussians,
Hard-Concrete gates, additive first-order effect composition, pairwise relations, kernel bandwidths, grouped task demonstrations
supplied by the dataset, gradient optimisation, and thresholded TP competence. Their role is to make the first test finite and
interpretable. Failure of one such choice does not by itself falsify P1--P8; failure across credible implementations and appropriate
coverage does.

\begin{example}
If the first model fails, that tells us the \emph{implementation} was wrong --- not necessarily that "the body has no reusable
mechanisms." We only conclude the prior is false if many different, reasonable implementations all fail under enough data.
\end{example}

\section*{4. Mathematical Modeling}

\begin{plainword}
This is the longest section. Each block of equations is followed by an italic plain-language translation. You can read only the
italic lines and the examples and still follow the whole argument.
\end{plainword}

\subsection*{4.1 Data, notation, and the two learning branches}

World-discovery data are temporally aligned interaction streams:

\[
\mathcal D_{\mathrm{world}}=
\left\{(O_0,a_0,\ldots,a_{T-1},O_T)\right\}.
\]

No object category, semantic action, TP boundary, $Z_D$, $R_D$, mechanism, or causal-graph label is required for training.
Simulation ground truth may be retained for evaluation only.

Task data provide episode boundaries and same-task grouping:

\[
\boxed{
\mathcal D_{\mathrm{task}}=\{\mathcal T_g\}_{g=1}^{G},\qquad
\mathcal T_g=\{\tau_{g,1},\ldots,\tau_{g,R_g}\},\quad R_g\geq2.}
\]

\begin{plainword}
For world learning we get raw "observe--act--observe" streams. For task learning we additionally get told which demonstrations
belong to the same task (but not \emph{what} the task is). Each task group has at least two demonstrations.
\end{plainword}

This grouping is dataset supervision for the first theory test, not a claim that task equivalence is already solved. The branches
share $S$, TP, and $W_D$:

\[
\boxed{
\mathcal D_{\mathrm{world}}\rightarrow S\rightarrow TP\rightarrow W_D,
\qquad
\mathcal D_{\mathrm{task}}\rightarrow TP\text{ sequences}\rightarrow C_T.}
\]

\subsection*{4.2 Interaction representation}

\subsubsection*{Temporally persistent latent state}

Visual tokens are

\[
H_t=E_{\mathrm{vis}}(O_t)=\{h_t^1,\ldots,h_t^N\}.
\]

The latent state is a bounded-capacity set

\[
\boxed{S_t=E_S(H_t,S_{t-1})=\{s_t^1,\ldots,s_t^{K_{\max}}\},}
\]

with optional activity variables. The executed action is \emph{not} input to $E_S$; state inference and intervention remain
distinct. The first implementation uses a ViT-style token encoder and a SAVi-like recurrent Slot Attention module. ``Slot'' denotes
a persistent latent interaction unit.

\begin{plainword}
The raw pixels are turned into a small fixed set of "slots" --- persistent placeholders that follow the same thing across time. The
action is deliberately kept out of this step, so that "what the state is" and "what we did" stay cleanly separate.
\end{plainword}

\begin{example}
Think of the slots as the rows of a chart: one row keeps tracking "the blood pressure," another "the heart rate," across time. The
row for blood pressure stays the same row as it changes, so we can talk about how \emph{much} it changed.
\end{example}

\subsubsection*{Action-conditioned dynamics and temporal correspondence}

The transition model is

\[
\boxed{\widehat S_{t+1}=F_S(S_t,a_t).}
\]

The low-level control should be the command actually executed, for example

\[
a_t=(\Delta x_t,\Delta y_t,\Delta z_t,\Delta r_t,g_t),
\]

rather than a semantic label such as ``push''. A small MLP embeds $a_t$; a permutation-equivariant slot Transformer processes
$S_t$; and FiLM injects the action condition:

\[
h_t^a=E_a(a_t),\qquad
(\gamma_t,\beta_t)=G_a(h_t^a),\qquad
\widetilde s_t^k=\gamma_t\odot \bar s_t^k+\beta_t.
\]

\begin{plainword}
The model predicts the next state from the current state plus the action. The action is a raw number (how far the gripper moved),
not a name like "push," because the model must learn from numbers what the action \emph{does}.
\end{plainword}

Residual slot-order ambiguity is resolved during training by

\[
\pi_t^*=\argmin_{\pi}\sum_{k=1}^{K_{\max}}
D_s(\widehat s_{t+1}^k,s_{t+1}^{\pi(k)}),
\]

yielding

\[
\boxed{L_{\mathrm{dyn}}=\sum_{t,k}
D_s(\widehat s_{t+1}^k,s_{t+1}^{\pi_t^*(k)}).}
\]

\begin{plainword}
The slots are not labelled, so the model may store them in any order. Before comparing predicted vs.\ actual, we find the best
matching of slots (a "Hungarian matching") so we compare like with like. The loss is the mismatch after this best matching.
\end{plainword}

The matched change is

\[
\Delta s_t^k=s_{t+1}^{\pi_t^*(k)}-s_t^k,
\qquad
\Delta S_t=\{\Delta s_t^k\}_{k=1}^{K_{\max}}.
\]

The minimal interaction-representation objective is

\[
\boxed{L_{\mathrm{IR}}=L_{\mathrm{dyn}}.}
\]

A weak observation reconstruction term $\lambda_{\mathrm{rec}}L_{\mathrm{rec}}$ may be used only to diagnose or prevent collapse.
Strong reconstruction is avoided because it can preserve appearance detail irrelevant to intervention-induced change.

\subsection*{4.3 Transformation primitives}

\subsubsection*{Experience and variable duration}

For segment $i$ with boundaries $b_i<e_i$,

\[
\boxed{
\mathcal I_i=(S_i^{\mathrm{pre}},A_i,S_i^{\mathrm{post}}),\quad
S_i^{\mathrm{pre}}=S_{b_i},\quad
S_i^{\mathrm{post}}=S_{e_i},\quad
A_i=a_{b_i:e_i-1}.}
\]

After matching, $\Delta S_i=S_i^{\mathrm{post}}-S_i^{\mathrm{pre}}$. A candidate duration satisfies $d_t\in\{1,\ldots,H_{\max}\}$.
Transition tokens

\[
x_\tau=E_b(S_\tau,a_\tau,\Delta S_\tau)
\]

are processed by a temporal Transformer. A hazard head produces $b_{t+j}=\sigma(H_b(h_{t+j}))$ and

\[
\boxed{P(d_t=j)=b_{t+j}\prod_{r<j}(1-b_{t+r}).}
\]

The boundary is meaningful only if one TP remains sufficient inside the segment and ceases to be sufficient beyond it.

\begin{plainword}
A "transformation" is a stretch of time over which a single description of "what changed" is enough. The model learns where these
stretches begin and end. The boundary is real only if one description works inside the stretch but fails once you cross the line.
\end{plainword}

\begin{example}
"Blood pressure fell steadily for 20 minutes" is one transformation. If at minute 21 it suddenly starts rising, that is a new
transformation --- the old description no longer fits. The model learns to cut the timeline at exactly such points.
\end{example}

\subsubsection*{Posterior: what happened}

The transformation posterior excludes action:

\[
\boxed{q_\phi(p_i\mid S_i^{\mathrm{pre}},S_i^{\mathrm{post}}).}
\]

Matched slot pairs form tokens

\[
x_i^k=E_\Delta(s_{i,\mathrm{pre}}^k,\Delta s_i^k).
\]

A Set Transformer and TP query pool the multi-slot change:

\[
H_i^\Delta=\mathrm{SetTransformer}_{\mathrm{post}}(\{x_i^k\}),
\qquad
h_i^{\mathrm{post}}=\mathrm{CrossAttn}(q_{\mathrm{TP}},H_i^\Delta).
\]

The first posterior is diagonal Gaussian,

\[
q_\phi=\mathcal N(\mu_i^q,\operatorname{diag}((\sigma_i^q)^2)),
\qquad
p_i=\mu_i^q+\sigma_i^q\odot\epsilon.
\]

Gaussianity is an implementation baseline, not a world prior.

\begin{plainword}
Given the state before and the state after, the "posterior" answers: \emph{what happened?} Crucially it is \emph{not} given the
action, so it must describe the change itself, not parrot back "you pushed."
\end{plainword}

\subsubsection*{Prior: what an intervention is expected to cause}

The prospective TP prior is

\[
\boxed{p_\theta(p_i\mid S_i^{\mathrm{pre}},A_i).}
\]

A Transformer encodes the variable-length action sequence, a set encoder summarises the pre-state, and a fusion network predicts
$(\mu_i^p,\log\sigma_i^p)$:

\[
h_i^A=E_A(A_i),\quad h_i^S=E_{\mathrm{set}}(S_i^{\mathrm{pre}}),\quad
h_i^{\mathrm{prior}}=F_{\mathrm{prior}}(h_i^S,h_i^A).
\]

The prior is also diagonal Gaussian in the first implementation. Posterior and prior are aligned by

\[
\boxed{
L_{\mathrm{align}}=
D_{\mathrm{KL}}\!\left[
q_\phi(p_i\mid S_i^{\mathrm{pre}},S_i^{\mathrm{post}})
\Vert
p_\theta(p_i\mid S_i^{\mathrm{pre}},A_i)
\right].}
\]

\begin{plainword}
The "prior" answers the forward question: \emph{given the state and the action, what change do we expect?} We train the two views to
agree. This single agreement objective enforces the two contrasts that give the TP its meaning (below).
\end{plainword}

This objective permits the two contrasts essential to the intended semantics:

\[
\begin{aligned}
S^{(1)}\neq S^{(2)},\ A^{(1)}\approx A^{(2)}
&\Rightarrow p^{(1)}\not\approx p^{(2)}\ \text{when effects differ},\\
A^{(1)}\neq A^{(2)},\ \Delta S^{(1)}\approx\Delta S^{(2)}
&\Rightarrow p^{(1)}\approx p^{(2)}.
\end{aligned}
\]

\begin{plainword}
(1) Same action, different states, different effects $\Rightarrow$ the TP must differ (it tracks the \emph{effect}, not the action).
(2) Different actions, same effect $\Rightarrow$ the TP must be the same (it merges equivalent outcomes).
\end{plainword}

\begin{example}
Contrast (1): the same dose of diuretic in a dehydrated vs.\ overloaded patient gives different outcomes, so the "what happened"
description differs. Contrast (2): lowering blood pressure by a diuretic vs.\ by a beta-blocker reaches the same outcome, so the
"what happened" description should merge. This is exactly what separates "transformation" from "action name."
\end{example}

\subsubsection*{TP effect and boundary semantics}

The TP must explain latent state change:

\[
\widehat{\Delta S}_i^{TP}=F_{\mathrm{TP}}(S_i^{\mathrm{pre}},p_i),
\qquad
\boxed{L_{\mathrm{effect}}^{TP}=D_S(\widehat{\Delta S}_i^{TP},\Delta S_i).}
\]

The first decoder is a TP-conditioned set-equivariant slot Transformer. For each time inside a segment,

\[
\widehat{\Delta S}_{i,t}^{TP}=F_{\mathrm{TP}}(S_t,p_i),
\]

and

\[
\boxed{L_{\mathrm{seg}}=\sum_i\sum_{t=b_i}^{e_i-1}
D_S(\widehat{\Delta S}_{i,t}^{TP},\Delta S_t).}
\]

A weak minimum-description-length pressure prevents a boundary at every timestep:

\[
L_{\mathrm{complexity}}=N_{\mathrm{segments}}
\quad\text{or}\quad
L_{\mathrm{complexity}}=\sum_t b_t.
\]

The TP objective is

\[
\boxed{
L_{\mathrm{TP}}=L_{\mathrm{effect}}^{TP}
+\lambda_A L_{\mathrm{align}}
+\lambda_S L_{\mathrm{seg}}
+\lambda_C L_{\mathrm{complexity}}.}
\]

\begin{plainword}
The TP must actually reproduce the observed change (effect), agree between its two views (align), fit every step inside its own
segment (seg), and avoid cutting the timeline into uselessly tiny pieces (complexity). These four pressures together define a good
transformation.
\end{plainword}

During this stage $L_{\mathrm{TP}}$ may update $E_S$, allowing interaction units and transformation coordinates to co-adapt. TP
dimensionality is selected by held-out prediction, reuse, and cross-context transfer; compactness is not assumed for aesthetic
reasons.

\subsection*{4.4 From TPs to the relational latent world $W_D$}

\subsubsection*{TP experience and many-to-many mechanism allocation}

Every TP retains its provenance:

\[
\boxed{
\mathcal E_i=(S_i^{\mathrm{pre}},A_i,S_i^{\mathrm{post}},\Delta S_i,p_i).}
\]

Let $M_{\max}$ be an overcomplete capacity, not the assumed true number of mechanisms:

\[
Z_D=\{Z_D^1,\ldots,Z_D^{M_{\max}}\}.
\]

For experience $i$, $m_i^j\in\{0,1\}$ indicates participation, $v_i^j\in\mathbb R^{d_v}$ is the realisation conditional on
participation, and

\[
\boxed{z_i^j=m_i^jv_i^j.}
\]

The distinction is necessary because $v_i^j=0$ cannot otherwise distinguish inactivity from an active zero-valued realisation. The
abstraction posterior is

\[
\boxed{q_\eta(m_i,v_i\mid p_i,S_i^{\mathrm{pre}}).}
\]

\begin{plainword}
We keep more candidate mechanisms ($M_{\max}$) than we expect to need. For each experience, each mechanism either participates or
not (a switch $m$), and if it participates it has a value $v$. We keep the switch and the value separate, so that "off" is not
confused with "on but currently zero."
\end{plainword}

\begin{example}
A patient's chart: each mechanism is a row. "ACE-inhibitor active? yes/no" is the switch; "how much effect now" is the value. If the
drug is off, that is different from the drug being on but having no current effect. The switch-vs-value split records this
difference.
\end{example}

Learned mechanism queries $q_D^j$ attend to an embedded TP and pre-state slots:

\[
X_i=\{E_p(p_i),E_s(s_{i,\mathrm{pre}}^1),\ldots,E_s(s_{i,\mathrm{pre}}^{K_{\max}})\},
\]

\[
\boxed{h_i^j=\mathrm{CrossAttn}(q_D^j,X_i).}
\]

Independent heads predict participation and realisation. The first implementation uses

\[
m_i^j\sim\mathrm{HardConcrete}(H_m(h_i^j)),
\]

and

\[
v_i^j=\mu_{v,i}^j+\sigma_{v,i}^j\odot\epsilon,
\qquad
(\mu_{v,i}^j,\log\sigma_{v,i}^j)=H_v(h_i^j).
\]

There is no categorical competition because several mechanisms may participate in one TP.

\begin{plainword}
Each mechanism is a learned "query" that looks at the change and the state and decides: am I involved, and with what value? Multiple
mechanisms can be involved at once --- there is no single winner.
\end{plainword}

\subsubsection*{Reusable mechanism functions}

Each candidate variable owns a reusable mechanism identity $\mathcal M_j$. Its effect in experience $i$ is

\[
\boxed{e_i^j=m_i^j\mathcal M_j(S_i^{\mathrm{pre}},v_i^j),}
\]

where $e_i^j$ lies in slot-change space. The first network is a shared set-dynamics Transformer with mechanism-specific query
conditioning and small adapters. A shared backbone supplies generic state processing; the query, realisation, and adapter preserve
mechanism identity.

Before relation learning, effects compose additively:

\[
\boxed{\widehat{\Delta S}_i^Z=\sum_{j=1}^{M_{\max}}e_i^j,}
\qquad
\boxed{L_{\mathrm{effect}}^D=\sum_iD_S(\Delta S_i,\widehat{\Delta S}_i^Z).}
\]

\begin{plainword}
For now, the total predicted change is simply the \emph{sum} of the active mechanisms' effects. This additive rule is a deliberate
starting point, not a claim that the world is additive. It stops a clever network from secretly doing the relation-work that we want
to study explicitly in the next step.
\end{plainword}

\begin{example}
Predicted change = (effect of drug A) + (effect of drug B). The sum is the baseline. Any part of the outcome that the sum cannot
explain (say A boosts B) is precisely what the \emph{relation} layer will be asked to capture.
\end{example}

An optional low-capacity decoder reconstructs the TP from active mechanism realisations:

\[
y_i^j=m_i^jE_Z(q_D^j,v_i^j),\quad
h_i^P=\mathrm{CrossAttn}(q_P,\{y_i^j\}),\quad
\widehat p_i=G_P(h_i^P),
\]

\[
\boxed{L_{\mathrm{TP-rec}}=D_P(p_i,\widehat p_i).}
\]

This is a weak hierarchical consistency term, not the primary grounding objective. Sparse local participation is encouraged by

\[
\boxed{L_{\mathrm{participation}}=
\sum_{i,j}\mathbb E[\mathbf 1(m_i^j\neq0)].}
\]

The Phase-I objective is

\[
\boxed{
L_{Z_D}^{(0)}=L_{\mathrm{effect}}^D
+\lambda_P L_{\mathrm{participation}}
+\lambda_T L_{\mathrm{TP-rec}},}
\]

with $\lambda_T$ weak or zero. A useful mechanism is not one that is merely frequent; it must use the same $\mathcal M_j$ to predict
held-out effects across diverse contexts and TP combinations.

\begin{plainword}
Three pressures: (1) the mechanisms must jointly reproduce the observed change; (2) only a few should be active at once (sparsity);
(3) weakly, the mechanisms should still reconstruct the original transformation. A mechanism earns its keep only if it predicts
\emph{new} situations with the same function.
\end{plainword}

\subsubsection*{Preserving the TP mediation}

Prospective mechanism engagement must pass through the TP prior:

\[
\boxed{
\widetilde q_\eta(m,v\mid S,A)=
\int q_\eta(m,v\mid p,S)p_\theta(p\mid S,A)\,dp.}
\]

Monte Carlo estimation draws $p^{(l)}\sim p_\theta(p\mid S,A)$ and then $(m^{(l)},v^{(l)})\sim q_\eta(\cdot\mid p^{(l)},S)$.
Expected engagement is

\[
\boxed{\rho_j(S,A)\approx\frac1L\sum_{l=1}^Lm_j^{(l)}.}
\]

A direct shortcut $(S,A)\rightarrow Z_D$ is excluded, because it would allow $Z_D$ to become a second action embedding and would
sever the proposed explanatory hierarchy.

\begin{plainword}
To predict which mechanism will be engaged by a future action, the model must go \emph{through} the transformation: action $\to$
expected change $\to$ engaged mechanisms. It is forbidden to jump straight from action to mechanism, because that shortcut would let
the mechanism layer just memorise actions instead of modelling change.
\end{plainword}

\subsubsection*{Directed relations and the corrected gate placement}

The domain relation set is

\[
\boxed{R_D=\{R_D^{j\rightarrow k}\}_{j\neq k}.}
\]

$R_D^{j\rightarrow k}$ is a reusable operator by which source mechanism $j$ modifies the effect expressed by target mechanism $k$.
It is neither co-occurrence nor a separate effect generator. Let

\[
\mathcal T_{j\rightarrow k}
(e_i^k;S_i^{\mathrm{pre}},v_i^j,v_i^k)
\]

be a directed modulation. The instance relation residual is defined \emph{without} a graph-existence gate:

\[
\boxed{
r_i^{j\rightarrow k}=
m_i^jm_i^k
\left[
\mathcal T_{j\rightarrow k}
(e_i^k;S_i^{\mathrm{pre}},v_i^j,v_i^k)-e_i^k
\right].}
\]

This notation separates three facts: mechanisms participate in an experience through $m_i^jm_i^k$; the relation operator has an
experience-specific realisation $r_i^{j\rightarrow k}$; and the candidate edge belongs to the learned domain structure according to a
global gate $g_{jk}$.

The first relation network is a shared directed residual-FiLM network with a pair embedding $q_R^{j\rightarrow k}$. With
$c_i=\mathrm{AttnPool}(S_i^{\mathrm{pre}})$,

\[
(\gamma_i^{jk},\beta_i^{jk})=
F_R(v_i^j,v_i^k,c_i,q_R^{j\rightarrow k}),
\]

\[
\widetilde e_{i,l}^{j\rightarrow k}
=(1+\gamma_i^{jk})\odot e_{i,l}^k+\beta_i^{jk},
\]

and hence

\[
r_i^{j\rightarrow k}=m_i^jm_i^k
(\widetilde e_i^{j\rightarrow k}-e_i^k).
\]

Candidate relation existence is a domain-global structural parameter

\[
\boxed{g_{jk}\sim\mathrm{HardConcrete}(\alpha_{jk}).}
\]

The gate is applied only when the world aggregates relation realisations:

\[
\boxed{
\widehat{\Delta S}_i^W=
\sum_j e_i^j+
\sum_{j\neq k}g_{jk}r_i^{j\rightarrow k}.}
\]

Equivalently,

\[
\widetilde e_i^k=e_i^k+\sum_{j\neq k}g_{jk}r_i^{j\rightarrow k},
\qquad
\widehat{\Delta S}_i^W=\sum_k\widetilde e_i^k.
\]

This corrected placement avoids defining the same structural gate both inside the relation realisation and again in the world graph.
The world grounding loss is

\[
\boxed{L_{\mathrm{world}}=\sum_iD_S(\Delta S_i,\widehat{\Delta S}_i^W).}
\]

\begin{plainword}
A relation $j\to k$ says: "when mechanism $j$ and mechanism $k$ are both active, $j$ changes the effect of $k$." The final predicted
change is the sum of all mechanism effects \emph{plus}, for each active relation, a small correction. The correction is gated by a
global switch $g_{jk}$ that says whether this relation is part of the learned world at all.
\end{plainword}

\begin{example}
Drug A and drug B each have their own effect (the sum). But A also \emph{potentiates} B --- B's effect is bigger when A is present.
The relation $A\to B$ captures exactly that extra "bigger than the sum" part. If the model finds this relation never helps, the gate
$g_{AB}$ turns it off.
\end{example}

The first implementation represents pairwise effect modulation observable when both mechanisms participate. Activation of an
otherwise absent target and stable higher-order interactions are outside this first-order model; they should be introduced only if
sufficiently supported residuals persist.

\subsubsection*{Directional relation identifiability}

Co-participation does not identify direction. Define the participation-aware source state

\[
\boxed{\bar z_i^j=(m_i^j,m_i^jv_i^j).}
\]

To support $j\rightarrow k$, experiences $a,b$ should approximately match pre-state and target realisation while source $j$ varies.
Let

\[
w_{ab}^{(k)}=m_a^km_b^k
K_S(S_a^{\mathrm{pre}},S_b^{\mathrm{pre}})
K_v(v_a^k,v_b^k),
\]

\[
d_{ab}^{(j)}=D_z(\bar z_a^j,\bar z_b^j).
\]

Directional contrastive support is

\[
\boxed{
\mathcal E_{j\rightarrow k}^{R}=
\Phi\!\left(\sum_{a<b}w_{ab}^{(k)}d_{ab}^{(j)}\right),}
\]

where $\Phi$ is monotone and saturating. This is an evidence statistic, not relation strength, and generally

\[
\mathcal E_{j\rightarrow k}^{R}\neq\mathcal E_{k\rightarrow j}^{R}.
\]

\begin{plainword}
Just seeing A and B together does not tell us \emph{which way} the arrow points. To support "A affects B," we look for pairs of
experiences where B's context is held fixed while A varies --- if the outcome changes with A, that is evidence for $A\to B$. The
reverse direction needs its own, different evidence.
\end{plainword}

\begin{example}
If A and B always appear together, we cannot tell whether A changes B or B changes A. We need cases where B is held constant and
A varies. In medicine this is exactly the idea behind a controlled experiment: vary one factor, hold the rest fixed, watch the
outcome.
\end{example}

To test whether an edge has functional value, let $\widehat{\Delta S}_i^{-jk}$ remove $g_{jk}r_i^{j\rightarrow k}$ from the world
prediction. The held-out explanatory power is

\[
\boxed{
EP_{j\rightarrow k}=
\mathbb E_{i\sim\mathcal D_{\mathrm{held}}}
\left[
D_S(\Delta S_i,\widehat{\Delta S}_i^{-jk})
-D_S(\Delta S_i,\widehat{\Delta S}_i^W)
\right].}
\]

\begin{plainword}
An edge earns its keep if, on data the model has never seen, removing the edge makes the prediction worse. That gap is its
explanatory power.
\end{plainword}

An edge is unresolved when $\mathcal E_{j\rightarrow k}^{R}$ is insufficient; it is supported only when evidence is sufficient and
$EP_{j\rightarrow k}$ remains positive across held-out contexts. Lack of evidence is not evidence of absence. Accordingly, let

\[
s_{jk}^R=f(\mathcal E_{j\rightarrow k}^R),
\]

with $s_{jk}^R\approx0$ under insufficient support. The support-conditioned edge penalty is

\[
\boxed{L_{\mathrm{edge}}^{\mathrm{support}}=
\sum_{j\neq k}s_{jk}^R\mathbb E[g_{jk}].}
\]

Only candidates that the data could have adjudicated pay a strong absence pressure.

\begin{plainword}
We only strongly penalise an edge for being "absent" when the data actually had a chance to prove it. If the data never varied the
right factor, we do not punish the edge for being uncertain --- "no evidence yet" is not the same as "no effect."
\end{plainword}

\begin{example}
If we never tested drug A in the presence of drug B, we cannot conclude "A does not affect B." We simply have no data. The penalty
must respect this: it only applies when the experiment \emph{could have} shown the effect.
\end{example}

\subsubsection*{$Z_D\leftrightarrow R_D$ co-refinement and staged training}

Let $\Sigma_j^{\mathrm{int}}$ summarise how interventions engage mechanism $j$, and define its relational signature

\[
\mathcal R_j=\{R_D^{j\rightarrow k},R_D^{k\rightarrow j}\}_{k\neq j}.
\]

Operational identity is

\[
\boxed{
\operatorname{Identity}(Z_D^j)=
(\mathcal M_j,\Sigma_j^{\mathrm{int}},\mathcal R_j).}
\]

Relations refine variable identity but do not wholly define it. The overcomplete bank permits soft splitting when an inactive query
acquires stable participation and soft deactivation when a redundant query's participation tends to zero.

Training proceeds in three phases:

\begin{enumerate}
\item \textbf{Candidate mechanisms.} Set $R_D=0$ and optimise $L_{Z_D}^{(0)}$. Mechanisms must acquire baseline reusable roles
before relation modules receive capacity.
\item \textbf{Relation residuals.} Freeze or strongly slow $q_\eta$ and $\mathcal M_j$; enable $F_R$ and $g_{jk}$; optimise
\[
\boxed{L_R=L_{\mathrm{world}}+\lambda_RL_{\mathrm{edge}}^{\mathrm{support}}.}
\]
Relations must explain residual structure rather than replacing mechanism baselines.
\item \textbf{Joint co-refinement.} Unfreeze both levels at smaller learning rates. If structurally single-mechanism experiences exist,
\[
\mathcal D_{\mathrm{base}}=\{i:\|m_i\|_0\approx1\},
\]
use
\[
\boxed{L_{\mathrm{base}}=
\mathbb E_{i\in\mathcal D_{\mathrm{base}}}
D_S\!\left(\Delta S_i,\sum_je_i^j\right).}
\]
If such data do not exist, set $\lambda_B=0$ rather than fabricating a baseline.
\end{enumerate}

\begin{plainword}
We learn in stages, like teaching in order: first learn the mechanisms alone, then learn the relations as \emph{corrections} on top,
then fine-tune both together slowly. This stops the relations from quietly doing the mechanisms' job.
\end{plainword}

The co-refinement objective is

\[
\boxed{
L_{W_D}=L_{\mathrm{world}}
+\lambda_RL_{\mathrm{edge}}^{\mathrm{support}}
+\lambda_PL_{\mathrm{participation}}
+\lambda_BL_{\mathrm{base}}
+\lambda_TL_{\mathrm{TP-rec}}.}
\]

\subsection*{4.5 Active exploration: $W_D$ back to interaction}

\subsubsection*{Structural evidence, not a learned $U_D$}

Uncertainty is not introduced as an independent latent $U_D$. The relevant question is how much evidence the current dataset
contains for mechanism reuse and directed relation identification. Retain reliable $Z_D$ experiences

\[
\mathcal D_t^Z=\{x_i^Z\}_{i=1}^{N_t},
\qquad
x_i^Z=(S_i^{\mathrm{pre}},m_i,v_i).
\]

For mechanism $j$, define

\[
\mathcal D_j=\{(S_i^{\mathrm{pre}},v_i^j):m_i^j\approx1\}.
\]

The first explicit evidence estimator uses

\[
K_j[a,b]=m_a^jm_b^j
K_S(S_a^{\mathrm{pre}},S_b^{\mathrm{pre}})
K_v(v_a^j,v_b^j),
\]

\[
\boxed{\mathcal E_j^Z=\log\det(I+\alpha_ZK_j).}
\]

Near-duplicate experiences add little; new supported contexts or realisations expand evidence. Relation evidence reuses
$\mathcal E_{j\rightarrow k}^R$. The total statistic is

\[
\boxed{
\mathcal E_D(\mathcal D_t^Z;W_D)=
\sum_j\mathcal E_j^Z+
\lambda_E^R\sum_{j\neq k}\mathcal E_{j\rightarrow k}^R.}
\]

It is a functional of data and the current structural hypothesis, not a learned object.

\begin{plainword}
"Uncertainty" here is not a mysterious extra variable. It is simply a number that measures \emph{how much varied evidence we have}
for each mechanism and relation. Repeating the same experiment twice adds little; a genuinely new situation adds a lot. We do not
learn this number with a network --- we compute it directly from the data.
\end{plainword}

\begin{example}
Seeing "drug A lowers blood pressure in 100 identical patients" gives you one fact, not 100 facts. One genuinely new situation (a
different co-morbidity, a different dose) teaches you far more. The evidence statistic rewards the second kind of data.
\end{example}

\subsubsection*{Expected structural evidence gain}

For a candidate intervention sequence $A$ at state $S_t$,

\[
p\sim p_\theta(p\mid S_t,A),
\qquad
(m,v)\sim q_\eta(m,v\mid p,S_t),
\]

which predicts $x_A^Z=(S_t,m,v)$. The expected evidence increment is

\[
\boxed{
J_{\mathrm{evidence}}(A)=
\mathbb E_{p,m,v}
\left[
\mathcal E_D(\mathcal D_t^Z\cup\{x_A^Z\};W_D)
-\mathcal E_D(\mathcal D_t^Z;W_D)
\right].}
\]

This single term replaces separate learned variable-, relation-, pair-, and uncertainty policies. Exploration supplies evidence;
ordinary $Z_D\leftrightarrow R_D$ learning decides which structure is needed.

\begin{plainword}
To choose the next action, ask: \emph{if I do this, how much new evidence will I likely gain?} We estimate this by simulating the
action through our model, adding the predicted result to our data, and measuring the increase. The action that maximises expected
new evidence is the one we take.
\end{plainword}

\subsubsection*{TP competence and the abstraction frontier}

Evidence completion alone confirms the current ontology. To search for omitted structure, compare the lower-level TP effect with the
higher-level world explanation, but only where TP is locally reliable. Define

\[
\boxed{
\mathcal C_{TP}=\{(S,A):
\text{local prior--posterior error, TP effect error, and support are calibrated}\}.}
\]

Historical estimates use

\[
D_{TP}[q_\phi(p\mid S,S')\Vert p_\theta(p\mid S,A)],
\qquad
D_S(\Delta S,F_{TP}(S,p^{\mathrm{post}})),
\]

and neighbourhood support in $(S,A)$. No new competence network is required initially.

For a predicted TP,

\[
F_W(S,p)=
\mathbb E_{q_\eta(m,v\mid p,S)}
\left[
\sum_je^j+\sum_{j\neq k}g_{jk}r^{j\rightarrow k}
\right],
\]

\[
\boxed{F_D(S,p)=D_S(F_{TP}(S,p),F_W(S,p)).}
\]

The prospective abstraction-frontier utility is

\[
\boxed{
J_{AF}(A)=
\mathbb E_{p\sim p_\theta(p\mid S_t,A)}[F_D(S_t,p)],
\quad (S_t,A)\in\mathcal C_{TP}.}
\]

\begin{plainword}
A second signal tells the agent where its current world model is \emph{incomplete}. Where the low-level "what changed" prediction and
the high-level "mechanism explanation" disagree, something is missing. The agent seeks out exactly those disagreements --- but only
where its low-level model is trustworthy, so it does not chase its own noise.
\end{plainword}

\begin{example}
If the low-level model reliably says "blood pressure fell 20 points" but the mechanism model can only explain "fell 5 points," there
is a 15-point gap the mechanisms do not account for. That gap is a sign of a missing mechanism. The agent goes looking where such
gaps appear.
\end{example}

Large discrepancy does not prove missing structure; it proposes an experiment. After real execution, error attribution decides
whether TP or $W_D$ failed.

\subsubsection*{Minimal exploration objective and optimisation}

The complete knowledge objective has two terms:

\[
\boxed{J_{\mathrm{exp}}(A\mid S_t)=
J_{\mathrm{evidence}}(A)+\lambda_FJ_{AF}(A).}
\]

Subject to physical, safety, and competence constraints,

\[
\boxed{
A_t^*=\argmax_{A\in\mathcal A_{\mathrm{valid}}(S_t)\cap\mathcal C_{TP}(S_t)}
J_{\mathrm{exp}}(A\mid S_t).}
\]

\begin{plainword}
Exploration has exactly two goals: gather new evidence, and probe the frontier where the model is incomplete. The next action is the
one that best serves these two goals, within safety limits.
\end{plainword}

The reparameterised TP prior supplies $\partial p/\partial A$, and relaxed mechanism gates supply $\partial(m,v)/\partial p$. The
first optimiser is projected gradient ascent with multiple safe initialisations:

\[
\boxed{
A^{(r+1)}=
\Pi_{\mathcal A_{\mathrm{valid}}\cap\mathcal C_{TP}}
\left[A^{(r)}+\eta_A\nabla_AJ_{\mathrm{exp}}\right].}
\]

Discrete modes such as gripper open/close are enumerated; continuous components are optimised within each mode. Only a short prefix
is executed before observation and replanning. A sampling planner is an evaluation fallback when gradients are unreliable, not a
second exploration policy. Before any model is competent, safe broad interactions form $\mathcal D_{\mathrm{seed}}$.

\subsubsection*{Post-interaction error attribution}

After executing $A$, infer the posterior TP from the real outcome. Define

\[
\boxed{
E_{\mathrm{prior}}=
D_{TP}[q_\phi(p\mid S^{\mathrm{pre}},S^{\mathrm{post}})
\Vert p_\theta(p\mid S^{\mathrm{pre}},A)],}
\]

\[
\boxed{
E_{TP}=D_S(\Delta S_{\mathrm{real}},
F_{TP}(S^{\mathrm{pre}},p^{\mathrm{post}})),}
\]

\[
\boxed{
E_W=D_S(\Delta S_{\mathrm{real}},
F_W(S^{\mathrm{pre}},p^{\mathrm{post}})),
\qquad
G_W^{\mathrm{obs}}=E_W-E_{TP}.}
\]

\begin{plainword}
After acting, we do an "autopsy" to find which layer was wrong. If the expected and actual transformation disagree, the
action$\to$transformation link failed. If that agrees but the TP cannot reproduce the change, the transformation model failed. If
both are fine but the world model is off, the abstraction failed. Each failure is attributed to the right layer.
\end{plainword}

Large $E_{\mathrm{prior}}$ indicates intervention-to-TP failure; small $E_{\mathrm{prior}}$ but large $E_{TP}$ indicates a TP
representation/effect failure; small lower-level errors but large $E_W$ indicate a world-abstraction failure. An experience can first
update TP and later, once reliably represented, contribute to $W_D$.

\subsection*{4.6 Learning task constraint $C_T$ from grouped demonstrations}

\subsubsection*{TP sequence extraction}

The frozen interaction and TP models segment every task demonstration:

\[
\tau_{g,r}\longrightarrow
P_{g,r}=(p_1^{g,r},\ldots,p_{N_{g,r}}^{g,r}),
\]

with associated boundary states $S_0^{g,r},\ldots,S_{N_{g,r}}^{g,r}$. The first dataset supplies the task episode boundary and
grouping $g$. It does not supply TP labels or a task code.

\subsubsection*{Task encoder}

A causal or bidirectional sequence Transformer with positional encodings maps the complete support demonstration to

\[
\boxed{c_{g,r}=E_C(P_{g,r})\in\mathbb R^{d_C}.}
\]

When initial and terminal states carry information needed to disambiguate otherwise identical TP sequences, their set summaries may
be included as tokens. This is an implementation interface; the scientific test is whether $c_{g,r}$ transfers functionally to other
realisations of task $g$.

\begin{plainword}
We read one demonstration of a task and compress it into a short "task code" $C_T$. The real test is not whether this code can
reconstruct its own demonstration, but whether it can predict a \emph{different} demonstration of the same task.
\end{plainword}

\subsubsection*{Explicit $W_D$ interface}

The task predictor does not receive only a monolithic state embedding. It receives an explicit world-structure readout containing
the currently instantiated mechanism nodes and gated directed relations. For query step $i$,

\[
(m_i,v_i)\sim q_\eta(m,v\mid p_i,S_i),
\]

\[
\mathcal W_i=
\left(
\{q_D^j,m_i^j,v_i^j,e_i^j\}_j,
\{g_{jk},r_i^{j\rightarrow k}\}_{j\neq k}
\right).
\]

A one-layer set/graph encoder $E_W$ produces query features while preserving node and directed-edge identity:

\[
h_i^W=E_W(S_i,\mathcal W_i).
\]

This computation is an on-demand readout of $W_D$, not a persistent task graph $G_T$.

\begin{plainword}
The task model is not given a vague "situation number." It is given the actual learned structure: which mechanisms are active, what
they are doing, and which relations are in play. This keeps the task model honest --- it must reason through the explicit world, not
around it.
\end{plainword}

\subsubsection*{Multimodal task progression}

Let

\[
x_i^{g,s}=(S_i^{g,s},h_i^{W,g,s},p_i^{g,s},c_{g,r})
\]

for a query demonstration $s\neq r$. The next-TP distribution is a conditional Gaussian mixture:

\[
\boxed{
P_T(p_{i+1}^{g,s}\mid x_i^{g,s})=
\sum_{q=1}^{Q}\pi_q(x_i^{g,s})
\mathcal N\!\left(
p_{i+1}^{g,s};\mu_q(x_i^{g,s}),
\operatorname{diag}(\sigma_q^2(x_i^{g,s}))
\right).}
\]

The start distribution omits a previous TP:

\[
P_T(p_1^{g,s}\mid S_0^{g,s},h_0^{W,g,s},c_{g,r}),
\]

and a Bernoulli head predicts

\[
P_T(y_i^{\mathrm{stop}}\mid S_i^{g,s},h_i^{W,g,s},p_i^{g,s},c_{g,r}).
\]

Multimodality prevents the mean of two valid strategies from becoming an invalid next TP and reduces pressure to encode strategy
identity in $C_T$.

\begin{plainword}
"Given the goal, what happens next?" often has several valid answers. We predict a \emph{mixture} of possibilities rather than a
single average, and separately predict "should we stop now." The mixture keeps genuinely different next steps separate.
\end{plainword}

\begin{example}
For the goal "lower blood pressure," the next valid step could be "more diuretic" or "start a beta-blocker." A mixture model keeps
these as two options. A single average would be a meaningless blend. The stop signal answers "have we reached the target yet?"
\end{example}

\subsubsection*{Cross-demonstration objective}

The support code from demonstration $r$ must explain a different query demonstration $s$ of the same task:

\[
\boxed{
L_{\mathrm{start}}^{r\rightarrow s}=
-\log P_T(p_1^{g,s}\mid
S_0^{g,s},h_0^{W,g,s},c_{g,r}),}
\]

\[
\boxed{
L_{\mathrm{next}}^{r\rightarrow s}=
-\sum_{i=1}^{N_s-1}
\log P_T(p_{i+1}^{g,s}\mid
S_i^{g,s},h_i^{W,g,s},p_i^{g,s},c_{g,r}),}
\]

\[
\boxed{
L_{\mathrm{stop}}^{r\rightarrow s}=
-\sum_{i=0}^{N_s}
\log P_T(y_i^{\mathrm{stop},g,s}\mid
S_i^{g,s},h_i^{W,g,s},p_i^{g,s},c_{g,r}).}
\]

The task objective is

\[
\boxed{
L_{C_T}=\mathbb E_{g,r\neq s}
\left[
L_{\mathrm{start}}^{r\rightarrow s}
+L_{\mathrm{next}}^{r\rightarrow s}
+\lambda_{\mathrm{stop}}L_{\mathrm{stop}}^{r\rightarrow s}
\right].}
\]

There is no required Euclidean consistency penalty $\|c_{g,r}-c_{g,s}\|$. Task codes are equivalent when they induce the same
predictive function on held-out realisations. Self-reconstruction alone is not used because it admits a trajectory-identity solution.

\begin{plainword}
Training: take the code from demonstration $r$ and make it predict the start, the next steps, and the stopping of a \emph{different}
demonstration $s$. We do not force the two codes to be numerically equal --- only to make the same predictions. This is what stops
the code from just memorising one trajectory.
\end{plainword}

\subsubsection*{Task state and task transformation}

The minimal task state is

\[
\boxed{Z_T=(S_t,C_T).}
\]

The last TP $p_t$ may be supplied to the next-TP head as a transient order-one progress cue, but it is not a new persistent
component of $Z_T$. If an experiment shows that longer history is necessary, a progress memory can be added as a tested extension
rather than assumed in advance.

Task transformation is the conditional operator

\[
\boxed{
\mathcal T_T(W_D,Z_T;p_t)=
P_T(p_{\mathrm{next}},y^{\mathrm{stop}}\mid
S_t,W_D,C_T;p_t).}
\]

$C_T$ acts on $W_D$ by selecting, restricting, and reweighting its feasible transformation structure. The output remains a
distribution. No independent $D_T$ is needed to store one desired TP, and no explicit $G_T$ is needed to store a task-specific
subgraph before evidence shows that such a persistent object improves explanation or transfer.

\subsubsection*{$C_T$ identifiability and necessity diagnostics}

If $(S_i,p_i,W_D)$ already predicts the query future, the optimum may ignore $C_T$. This is not necessarily optimisation failure; the
dataset has not made task information necessary. Train a no-context baseline

\[
P_0(p_{\mathrm{next}},y^{\mathrm{stop}}\mid S,p,W_D)
\]

and report

\[
\boxed{
\Delta_{C_T}=
\mathbb E_{\mathrm{held}}
\left[
\log P_T(p_{\mathrm{next}},y^{\mathrm{stop}}\mid S,p,W_D,C_T)
-\log P_0(p_{\mathrm{next}},y^{\mathrm{stop}}\mid S,p,W_D)
\right].}
\]

Positive held-out $\Delta_{C_T}$ under new initial states and demonstrations is evidence that $C_T$ contributes information
unavailable from the current world state. Strategy invariance is separately tested by transferring $c_{g,r}$ across demonstrations
with different valid TP paths. World necessity is tested by comparing the full model against a predictor that removes $W_D$ under the
same task but different world configurations.

\begin{plainword}
We must check that the task code is actually \emph{doing something}. If a model without the task code predicts just as well, then the
task code is unnecessary --- the data did not require it. The gap between "with task code" and "without task code" on held-out data
is our measure of necessity.
\end{plainword}

\begin{example}
If the state alone ("BP is high") already tells you the next step, then "the goal is to lower BP" adds nothing. The task code earns
its place only when it changes the prediction in situations the state alone cannot disambiguate.
\end{example}

\subsection*{4.7 Task action generation by TP-prior inversion}

\subsubsection*{Distribution-level objective}

At execution time, infer $C_T$ from one or more support demonstrations, observe $S_t$, instantiate the explicit $W_D$ readout, and
obtain the desired distribution

\[
P_T^*(p)=P_T(p_{\mathrm{next}}=p\mid S_t,W_D,C_T;p_t).
\]

For candidate low-level intervention sequence $A$, the learned TP prior predicts

\[
P_A(p)=p_\theta(p\mid S_t,A).
\]

Action realisation is distribution matching:

\[
\boxed{
A_t^*=\argmin_{A\in\mathcal A_{\mathrm{valid}}\cap\mathcal C_{TP}}
\left\{
D_{TP}(P_A,P_T^*)+\lambda_aC_a(A)
\right\}.}
\]

where $C_a$ penalises unsafe, excessive, or unnecessarily long control. A Monte Carlo cross-entropy form that remains differentiable
through $p_\theta$ is

\[
\boxed{
D_{TP}(P_A,P_T^*)=
\mathbb E_{p\sim p_\theta(p\mid S_t,A)}[-\log P_T^*(p)].}
\]

To avoid exploiting an overly broad $P_A$, an entropy or variance penalty may be included as a controller regulariser. It is not a
task-model loss.

\begin{plainword}
To actually \emph{do} the task, we work backwards. The task model tells us the distribution of "good next changes." We then search
over actions for the one whose predicted change best matches that desired distribution. In plain terms: know what change you want,
then find the action that produces it.
\end{plainword}

\begin{example}
The goal says "lower blood pressure by a moderate amount." We do not have a memorised "lower-BP action." Instead we ask the world
model "which action would produce a moderate fall?" and pick that action. The action is found by inverting the model, not by a
separate policy.
\end{example}

\subsubsection*{First execution procedure}

The first implementation uses the same projected multi-start gradient optimiser as exploration, but minimises the task objective.
For highly separated TP modes, select or sample a valid mixture component $q^*$ and minimise expected distance to its mean under
component confidence. Enumerate discrete control modes, optimise continuous controls, execute a short prefix, observe, update $S_t$,
and replan. Terminate when the stop posterior and safety checks agree.

No neural task motor policy $\pi_{\omega}^{\mathrm{task}}$ and no extra supervised action loss are introduced. Exploration and task
execution reuse the same learned map from intervention to transformation:

\[
\boxed{
\begin{array}{rcl}
W_D&\rightarrow&J_{\mathrm{exp}}\rightarrow A,\\
(W_D,S_t,C_T)&\rightarrow&P_T^*(p)\rightarrow p_\theta(p\mid S_t,A)^{-1}\rightarrow A.
\end{array}}
\]

This is the unifying control principle of the framework: higher levels assign value to transformations; the TP prior supplies their
motor realisation.

\begin{plainword}
One and the same "action $\to$ change" model is used twice: forward for exploration, and inverted for task execution. The only
difference is what value the higher levels place on each possible change.
\end{plainword}

\subsection*{4.8 Gradient boundaries and staged optimisation}

The hierarchy is intended to abstract lower-level evidence, not rewrite it to fit a preferred high-level ontology. The first
implementation therefore uses explicit stop-gradient boundaries.

During interaction and TP learning,

\[
L_{\mathrm{IR}},L_{\mathrm{TP}}
\rightarrow E_{\mathrm{vis}},E_S,F_S,q_\phi,p_\theta,F_{TP}.
\]

During world abstraction,

\[
\boxed{\operatorname{stopgrad}(p),\qquad
L_{W_D}\rightarrow q_\eta,\mathcal M_j,F_R,g_{jk},}
\]

but $L_{W_D}\nrightarrow q_\phi,p_\theta,E_S$. During task learning,

\[
\boxed{
\operatorname{stopgrad}(TP),\quad
\operatorname{stopgrad}(W_D),\quad
L_{C_T}\rightarrow E_C,P_T,E_W,}
\]

where $E_W$ is only the task-side readout of the frozen explicit world structure. During exploration or task execution, all model
parameters are frozen and gradients update only $A$.

\begin{plainword}
"Stop-gradient" means: when the higher level is learning, the lower level's parameters are \emph{frozen}. The higher level must
explain what the lower level already produced, and is not allowed to reach down and reshape the lower level to make itself look good.
\end{plainword}

\begin{example}
A professor grading students should not be allowed to change the students' answers to make the grading look successful. Similarly,
the mechanism layer must explain the transformations as they are --- it cannot quietly edit the transformations to fit its theory.
\end{example}

\begin{center}
\small
\begin{tabularx}{\textwidth}{@{}lXX@{}}
\toprule
Stage & Optimised quantities & Frozen boundary \\
\midrule
Interaction/TP & state encoder, dynamics, TP posterior/prior/effect, boundary & none between $S$ and TP \\
Candidate $Z_D$ & allocation and mechanism bank & TP representation \\
Relation discovery & relation network and edge gates & TP; initially mechanism bank \\
Co-refinement & $Z_D$ and $R_D$ at small learning rates & TP \\
Exploration & candidate intervention $A$ & all model parameters \\
$C_T$ learning & task encoder, multimodal predictor, $W_D$ readout & TP and $W_D$ \\
Task execution & candidate intervention $A$ & all model parameters \\
\bottomrule
\end{tabularx}
\end{center}

High-level feedback still exists through data: $W_D$ changes exploration, exploration changes the interaction distribution, and new
reliable interactions update TP and $W_D$. This is system-level feedback rather than a direct gradient that could make lower evidence
conform to the current ontology.

\subsection*{4.9 Unified objectives and closed loops}

The training objectives are

\[
\boxed{L_{IR}=L_{\mathrm{dyn}},}
\]

\[
\boxed{L_{TP}=L_{\mathrm{effect}}^{TP}
+\lambda_AL_{\mathrm{align}}
+\lambda_SL_{\mathrm{seg}}
+\lambda_CL_{\mathrm{complexity}},}
\]

\[
\boxed{L_{Z_D}^{(0)}=L_{\mathrm{effect}}^D
+\lambda_PL_{\mathrm{participation}}
+\lambda_TL_{\mathrm{TP-rec}},}
\]

\[
\boxed{L_R=L_{\mathrm{world}}
+\lambda_RL_{\mathrm{edge}}^{\mathrm{support}},}
\]

\[
\boxed{L_{W_D}=L_{\mathrm{world}}
+\lambda_RL_{\mathrm{edge}}^{\mathrm{support}}
+\lambda_PL_{\mathrm{participation}}
+\lambda_BL_{\mathrm{base}}
+\lambda_TL_{\mathrm{TP-rec}},}
\]

and the cross-demonstration $L_{C_T}$ defined above. Exploration and task action selection are action-space objectives, not
model-training losses.

The world-discovery loop is

{\footnotesize
\[
\boxed{
Interaction\rightarrow TP\rightarrow Z_D\leftrightarrow R_D\rightarrow W_D
\rightarrow
\left\{
\begin{array}{c}
Evidence\ Completion\\
Abstraction\ Challenge
\end{array}
\right.
\rightarrow Intervention\rightarrow New\ Interaction.}
\]}

\begin{plainword}
World learning is a loop: interact $\to$ describe changes $\to$ build mechanisms and relations $\to$ form a world model $\to$ either
fill in missing evidence or probe the frontier $\to$ act $\to$ learn from the new outcome, and around again.
\end{plainword}

The task loop is

{\footnotesize
\[
\boxed{
\begin{aligned}
Grouped\ Demonstrations&\rightarrow TP\ Sequences\rightarrow C_T,\\
(S_t,C_T)+W_D&\rightarrow P_T(p_{\mathrm{next}},stop)
\rightarrow TP\ Prior\ Inversion\\
&\rightarrow A_t\rightarrow S_{t+1}\rightarrow Replan.
\end{aligned}}
\]}

\begin{plainword}
Task loop: several demonstrations of one task $\to$ a task code $\to$ (with the world model) a prediction of the next good change and
whether to stop $\to$ invert to find the action $\to$ act $\to$ observe $\to$ replan.
\end{plainword}

Together they yield the full theory:

\[
\boxed{
\begin{aligned}
Interaction\ Representation&\rightarrow TP
\rightarrow
\left\{
\begin{array}{l}
Z_D\leftrightarrow R_D\rightarrow W_D\rightarrow exploration,\\
TP\ sequences\rightarrow C_T
\end{array}
\right.\\
&\rightarrow task\ transformation\ on\ W_D
\rightarrow action.
\end{aligned}}
\]

\section*{5. Computation and Experiments --- Future Work}

\subsection*{5.1 Status of this section}

This section is a preregistered computational plan. It specifies the first implementation and the observations that would support or
falsify each layer. It reports no completed experiment. Its purpose is to prevent later empirical convenience from silently
redefining the theory.

\begin{plainword}
"Preregistered" means: we write down what would count as success or failure \emph{before} running the experiment, so we cannot move
the goalposts afterwards. This is the scientific equivalent of registering a clinical trial protocol before enrolling patients.
\end{plainword}

\subsection*{5.2 First-implementation network and algorithm specification}

\begin{center}
\small
\begin{tabularx}{\textwidth}{@{}p{0.18\textwidth}p{0.35\textwidth}X@{}}
\toprule
Component & First implementation & Reason for this first test \\
\midrule
$E_{\mathrm{vis}}$ & ViT patch-token encoder & localised token interface for latent grouping \\
$E_S$ & SAVi-like recurrent Slot Attention & persistent bounded-capacity set of interaction units \\
$F_S$ & set Transformer with action FiLM & state-dependent intervention effects without concatenating action into state inference \\
correspondence & Hungarian matching & resolve residual training-time slot permutation \\
TP boundary & temporal Transformer with hazard head & variable duration and explicit termination probability \\
$q_\phi(p\mid S^{pre},S^{post})$ & Set Transformer plus TP query; diagonal Gaussian & permutation-aware multi-slot effect encoding and reparameterised baseline \\
$p_\theta(p\mid S^{pre},A)$ & state-set encoder plus action Transformer; diagonal Gaussian & prospective intervention-to-transformation map \\
$F_{TP}$ & TP-conditioned set Transformer & distribute one TP effect across persistent interaction units \\
$q_\eta(m,v\mid p,S)$ & mechanism-query cross-attention & simultaneous many-to-many allocation \\
$m^j$ & Hard-Concrete participation gate & differentiable sparse local activation \\
$v^j$ & diagonal-Gaussian realisation head & context-specific stochastic realisation \\
$\mathcal M_j$ & shared set-dynamics backbone plus mechanism adapter & reuse generic dynamics while retaining mechanism identity \\
$R_D^{j\rightarrow k}$ & shared directed residual-FiLM network plus pair embedding & constrain an edge to modify a target mechanism effect \\
$g_{jk}$ & global Hard-Concrete edge gate & sparse domain-level candidate relation structure \\
world executor & one directed message-passing layer and explicit gated sum & preserve inspectable node/edge contribution \\
mechanism evidence & kernel log-determinant & explicit non-redundant coverage statistic \\
relation evidence & matched-context directional kernel support & measure whether source variation exists under comparable target context \\
exploration & projected multi-start gradient optimisation, receding horizon & use the differentiable TP bridge without a learned exploration policy \\
$E_C$ & sequence Transformer over TP tokens & encode long-horizon persistent task constraint \\
$P_T$ & task/world-conditioned Gaussian mixture plus stop head & represent multiple valid next transformations \\
task control & TP-prior distribution matching, projected multi-start optimisation, receding horizon & reuse intervention competence without a separate task motor policy \\
\bottomrule
\end{tabularx}
\end{center}

These choices define the first falsifiable system. If gradients are numerically unreliable, a sampling optimiser is a controller
ablation; if a diagonal-Gaussian TP cannot represent observed transformations, the distribution may be enriched. Such changes test
implementation adequacy and do not alter the physical priors unless failure persists across credible realisations.

\subsection*{5.3 Training-data design implied by the theory}

\subsubsection*{World data}

The unit of collection is the continuous, time-aligned trajectory, not an isolated image/action pair and not a video pre-cut at
semantic action boundaries. The first dataset is organised as

\[
\boxed{
\mathcal D_{\mathrm{world}}=
\mathcal D_{\mathrm{broad}}
\cup\mathcal D_{\mathrm{contrast}}
\cup\mathcal D_{\mathrm{explore}}.}
\]

\begin{itemize}
\item $\mathcal D_{\mathrm{broad}}$ uses safe low-level coverage to bootstrap $S$ and TP.
\item $\mathcal D_{\mathrm{contrast}}$ varies initial conditions and controllable physical configurations to expose diverse mechanism
combinations and directional contrasts. Environment design is not latent supervision: the model still receives only observations and
executed interventions.
\item $\mathcal D_{\mathrm{explore}}$ is collected by the learned evidence/frontier objective after TP calibration.
\end{itemize}

Five coverage axes matter more than raw sample count:

\begin{enumerate}
\item \textbf{State coverage:} comparable interventions in different $S$.
\item \textbf{Intervention coverage:} different $A$ from comparable $S$.
\item \textbf{Mechanism-combination coverage:} reusable mechanisms occur alone when possible and in varied combinations rather than
always co-varying.
\item \textbf{Realisation coverage:} each active mechanism spans multiple $v^j$ values and contexts.
\item \textbf{Directional contrast coverage:} target $k$ and its context are approximately controlled while source $j$ varies, and
separate data support the reverse direction.
\end{enumerate}

\begin{plainword}
What matters is not "more data" but "data with the right contrasts": the same action in different states, different actions in the
same state, each mechanism seen alone and in combinations, each mechanism seen at many strengths, and each relation tested in both
directions.
\end{plainword}

The first proof-of-concept should use a simulator in which the experimenter knows, but does not train on, controllable mechanisms
such as contact, friction, support, containment, attachment, collision, blocking, and grasp-mediated motion. Ground truth then
evaluates whether learned variables are reusable and structurally aligned without supplying their names.

\subsubsection*{Task data}

Every task group must contain multiple demonstrations. For the first experiment, $R_g=5$--$10$ is a practical target, although the
objective only requires $R_g\geq2$. Within a task group, demonstrations deliberately vary:

\begin{itemize}
\item initial state;
\item valid strategy and TP path;
\item world configuration, including constraints and obstacles;
\item low-level execution details such as speed and precise motor trajectory.
\end{itemize}

Across task groups, the dataset must contain locally similar states and TPs with different valid futures:

\[
\boxed{
(S_i^g,p_i^g,W_D)\approx(S_j^h,p_j^h,W_D),\ g\neq h,
\quad
P(p_{\mathrm{next}}\mid C_T^g)
\neq P(p_{\mathrm{next}}\mid C_T^h).}
\]

\begin{plainword}
For the task code to be learnable, there must be situations that look the same in the world but belong to different tasks, so the
only thing that decides "what next" is the task itself.
\end{plainword}

\begin{example}
"Blood pressure is high" looks the same, but the task might be "lower it" or "keep it steady for surgery." The world state alone
cannot decide; the task code must. That is why such data are necessary.
\end{example}

Otherwise $C_T$ is not identifiable as a necessary predictor. Within one task, the data must also contain different world
configurations requiring different valid next TPs:

\[
\boxed{
C_T\ \text{fixed},\quad (S,W_D)\ \text{varies}
\quad\Rightarrow\quad
P_T(p_{\mathrm{next}})\ \text{varies}.}
\]

Otherwise the task predictor can memorise a fixed script and ignore $W_D$. Cases with the same start and end but different middle
processes test $C_T\neq$ trajectory; cases where the start and end are similar but a specific intermediate transformation is
required test $C_T\neq$ terminal state.

Finally, task TP support must overlap the action-realisable world TP support:

\[
\boxed{
\operatorname{supp}(TP_{\mathrm{task}})
\subseteq\operatorname{supp}_{\mathrm{competent}}(p_\theta(p\mid S,A))
\ \text{approximately}.}
\]

This separates failure of task inference from failure to realise a correctly requested TP.

\begin{plainword}
The transformations the task demands must be ones the world model can actually produce with some action. Otherwise we cannot tell
"the task was understood wrong" from "the task was understood right but could not be performed."
\end{plainword}

\subsection*{5.4 Staged computational protocol}

\begin{enumerate}
\item Collect $\mathcal D_{\mathrm{seed}}$ with broad safe low-level controls and strict sensor/action synchronisation.
\item Train $E_{\mathrm{vis}},E_S,F_S$ and validate temporal state correspondence under held-out trajectories.
\item Train boundary inference, $q_\phi$, $p_\theta$, and $F_{TP}$; validate both action/state contrasts and action realisability.
\item Freeze TP; train candidate $Z_D$ without relations; select weak sparsity by held-out mechanism reuse, not by the prettiest
decomposition.
\item Freeze or slow $Z_D$; train directional relation residuals only where contrast support is nontrivial.
\item Jointly co-refine $Z_D\leftrightarrow R_D$ at small learning rates; report gauge sensitivity across seeds.
\item Activate evidence/frontier exploration; interleave short collection batches with recalibration and staged retraining.
\item Freeze TP and $W_D$; segment grouped task demonstrations and train $E_C,P_T$ with support/query pairs $r\neq s$.
\item Evaluate $C_T$ and $W_D$ necessity before attempting closed-loop control.
\item Execute task control by TP-prior inversion with receding-horizon replanning and error attribution.
\end{enumerate}

\begin{plainword}
This is the recipe, in order: gather data $\to$ learn state $\to$ learn transformations $\to$ learn mechanisms $\to$ learn relations
$\to$ refine both $\to$ explore actively $\to$ learn the task $\to$ verify the task and world are necessary $\to$ act. Each step is
validated before moving on.
\end{plainword}

\subsection*{5.5 Identifiability, necessity, and falsification diagnostics}

\begin{center}
\small
\begin{tabularx}{\textwidth}{@{}p{0.17\textwidth}XX@{}}
\toprule
Claim & Required diagnostic & Evidence against the claim \\
\midrule
Persistent $S$ & matching stability, held-out action-conditioned transition error, slot perturbation consistency & arbitrary rematching or equal performance from unstable framewise latents \\
TP is transformation, not action & same-action/different-state and different-action/similar-effect retrieval; prior--posterior calibration & TP predicted from action alone or unable to merge equivalent effects \\
TP boundary is meaningful & within-segment effect sufficiency and boundary perturbation test & boundaries coincide mechanically with action switches or every timestep \\
$Z_D$ is reusable & cross-context and held-out-combination prediction by fixed $\mathcal M_j$ & query-specific memorisation, universal participation, or no compositional transfer \\
$R_D$ is directed structure & $\mathcal E_{j\rightarrow k}^R$, edge removal $EP_{j\rightarrow k}$, reverse-direction comparison & edge survives without contrast, has no held-out gain, or flips across seeds/contexts \\
$W_D$ adds abstraction & world prediction relative to TP decoder, held-out recombination, description length & monolithic TP decoder matches all performance and structure is unstable \\
exploration is structural & evidence gained per interaction, edge/variable resolution rate, frontier discoveries & novelty/random coverage matches or beats the objective at equal budget \\
$C_T$ is necessary & held-out $\Delta_{C_T}$ on locally matched cross-task states & no-context predictor performs equally \\
$C_T$ is strategy-invariant & cross-strategy $r\rightarrow s$ likelihood and task retrieval by functional transfer & code follows trajectory/demonstrator rather than task group \\
$C_T$ acts on $W_D$ & same-task/different-world next-TP prediction; remove-$W_D$ ablation & fixed script succeeds equally or $W_D$ removal has no cost \\
Task TP is realisable & distribution match before/after action, execution success within TP competence & correct next-TP prediction but systematic inversion failure \\
No $D_T/G_T$ needed & compare minimal operator with explicit-latent/graph ablations only after minimal failure & persistent residual requiring stable desired state or task graph \\
\bottomrule
\end{tabularx}
\end{center}

\begin{plainword}
For every claim the framework makes, this table states the experiment that would confirm it and the result that would refute it.
For example, "the TP is a transformation, not an action code" is confirmed if the same action in different states gives different TPs
and different actions giving the same effect give the same TP; it is refuted if the TP can be predicted from the action alone.
\end{plainword}

Latent alignment metrics against simulator ground truth are secondary. Because representations have permutation and
invertible-coordinate freedoms, the primary measurements are intervention prediction, held-out recombination, cross-context reuse,
relation removal, cross-demonstration prediction, and control consequences.

\subsection*{5.6 Core ablations}

The first study should isolate scientific necessity rather than enumerate every architecture. The mandatory ablations are:

\begin{enumerate}
\item action-conditioned versus action-free latent dynamics;
\item posterior without action versus posterior with action, to expose action-code leakage;
\item fixed one-step TP versus learned variable duration;
\item one-to-one TP clustering versus many-to-many mechanism allocation;
\item mechanisms only versus $Z_D+R_D$;
\item naive edge sparsity versus support-conditioned edge sparsity;
\item random/novelty exploration versus evidence only versus evidence plus abstraction frontier;
\item self-demonstration task reconstruction versus cross-demonstration prediction;
\item unimodal versus mixture next-TP prediction;
\item no $C_T$, full $C_T$, and shuffled task grouping;
\item no $W_D$, frozen explicit $W_D$, and a capacity-matched monolithic world embedding;
\item direct behaviour cloning versus TP-prior inversion at equal task data budget.
\end{enumerate}

\begin{plainword}
An "ablation" is the scientific version of "remove one part and see if it was doing anything." Each line removes or swaps one piece
and asks whether the result gets worse. This isolates which pieces are \emph{necessary} rather than decorative.
\end{plainword}

\subsection*{5.7 Failure modes to monitor}

\begin{itemize}
\item \textbf{Slot gauge failure:} temporally unstable interaction units make change meaningless.
\item \textbf{TP leakage:} the posterior memorises the post-state or the prior becomes an action embedding.
\item \textbf{Mechanism collapse:} all $m^j$ activate or one mechanism explains all effects.
\item \textbf{Mechanism/relation gauge exchange:} $\mathcal M_j$ and $R_D$ trade explanatory mass without stable held-out identity.
\item \textbf{False edge absence:} sparse penalties prune an edge before directional evidence exists.
\item \textbf{Ontology confirmation:} evidence-only exploration never proposes transformations outside the current $Z_D$ support.
\item \textbf{Model exploitation:} action optimisation finds adversarial controls outside calibrated TP support.
\item \textbf{Task-code shortcut:} $C_T$ encodes initial state, strategy, trajectory length, or demonstrator identity.
\item \textbf{World bypass:} $P_T$ predicts a memorised task script and ignores explicit $W_D$.
\item \textbf{Control confounding:} task inference is blamed for an unrealisable TP or a poor local optimiser.
\item \textbf{Partial observability:} hidden state makes identical observed $S$ respond differently; longer memory or additional
sensing is then required.
\end{itemize}

\begin{plainword}
These are the ways the project is most likely to fail, written down in advance so that a failure can be \emph{diagnosed} rather than
just noticed. Each failure has a known signature and a known fix.
\end{plainword}

\section*{6. Preliminary Experimental Results Analysis --- Future Work}

\subsection*{6.1 Status of this section}

No preliminary result is claimed in the present positioning paper. This section is intentionally reserved so that the eventual
analysis follows the causal order of the model rather than reporting only end-task success.

\subsection*{6.2 Required reporting order}

Results will be analysed in the following order:

\begin{enumerate}
\item data coverage and calibration, including the five world coverage axes and task support/query variation;
\item interaction-state stability and TP competence;
\item mechanism reuse and participation structure;
\item directional relation evidence and held-out explanatory power;
\item exploration efficiency and discoveries at the abstraction frontier;
\item $C_T$ necessity, cross-strategy transfer, and multimodal progression;
\item dependence of task prediction on explicit $W_D$;
\item TP-prior inversion accuracy and closed-loop task success;
\item robustness across seeds, domains, and representation gauges.
\end{enumerate}

An end-to-end success rate is interpretable only after the lower-level competence conditions have been reported. Otherwise a failure
cannot be assigned to perception, TP segmentation, intervention prediction, world abstraction, task inference, or action optimisation.

\begin{example}
In a clinical trial you would not report only "the patient recovered" --- you would report each link in the chain: was the diagnosis
right, was the drug given correctly, did the drug do what was expected, did the goal match the mechanism? The same discipline applies
here.
\end{example}

\subsection*{6.3 Result-interpretation template}

For every major claim, the eventual text should state:

\begin{enumerate}
\item the data contrast that made the variable or relation identifiable;
\item the held-out intervention or demonstration split;
\item the relevant baseline and capacity control;
\item uncertainty across random seeds;
\item whether improvement is predictive, structural, exploratory, or control-level;
\item which lower-level competence conditions held;
\item the narrowest warranted conclusion and the guarantees that remain absent.
\end{enumerate}

The strongest acceptable early claim is not that the model has discovered ``the true physics.'' It is that, in a controlled domain
with known but unlabelled generative mechanisms, the learned $Z_D/R_D$ system provides more stable cross-context effect decomposition
and more efficient evidence-directed interaction than plausible alternatives, while $C_T$ transfers across strategies and changes
transformation predictions beyond what $S$ and $W_D$ alone explain.

\section*{7. Research Position and Open Problems}

\subsection*{7.1 The central position}

The proposal takes a middle position between end-to-end control and fully symbolic world recovery. An agent does not need a complete
human ontology before acting, but it needs more than a predictive black box if it is to decide which uncertainty matters, recognise
when its ontology is incomplete, and reuse world knowledge across tasks. Transformation is the appropriate interface because it is
simultaneously observable after action, predictable before action, decomposable into reusable mechanisms, and invertible into control.

\begin{example}
The framework sits between two extremes. At one end, a pure black box that says "do this now" with no explanation (like a reflex).
At the other end, a full textbook of physics and anatomy written out in advance. This work claims the useful middle: a compact,
reusable model that is learned from experience and that can be \emph{inspected} and \emph{reused}, but that does not pretend to be a
complete theory.
\end{example}

The framework's strongest conceptual commitments are therefore:

\begin{enumerate}
\item effective variables should be defined by reusable effects, not visual axes alone;
\item relations should be identified by directional contrasts and held-out necessity, not co-occurrence;
\item uncertainty should remain an evidence status of a structural hypothesis, not automatically become another latent object;
\item task identity should be a cross-realisation predictive constraint, not a trajectory code;
\item task transformation should be the action of $C_T$ on explicit $W_D$, not an additional ontological layer unless experiments
require one;
\item exploration and task execution should reuse one intervention-to-transformation model and differ primarily in how higher levels
value transformations.
\end{enumerate}

\subsection*{7.2 Open scientific questions}

Several questions remain genuinely open rather than being disguised as design choices.

\begin{itemize}
\item Under what nonlinear conditions can TP-mediated effective variables be identifiable beyond permutation and invertible
reparameterisation?
\item Which intervention-coverage criterion is sufficient to separate mechanism baselines from directed relation residuals?
\item When does a pairwise relational world remain adequate, and how should higher-order structure be introduced without destroying
interpretability?
\item Can an abstraction frontier reliably expose unknown mechanisms when the lower TP model is itself trained on incomplete coverage?
\item What minimal diversity of demonstrations makes $C_T$ distinguish task constraint from strategy and terminal state?
\item When is an order-one TP progress cue sufficient, and when is explicit task-progress memory unavoidable?
\item Can one derive calibrated bounds linking TP-prior error, task distribution-matching error, and closed-loop task failure?
\item Across which domain changes does $W_D$ remain reusable, and when should the agent infer that it has entered a new domain rather
than refine the old one?
\end{itemize}

\subsection*{7.3 Final definition}

The complete research object is

\[
\boxed{
\begin{aligned}
O_t&\rightarrow S_t,\\
(S^{\mathrm{pre}},A,S^{\mathrm{post}})&\rightarrow p,\\
p&\rightarrow (m,v)\rightarrow Z_D,\\
Z_D&\leftrightarrow R_D\rightarrow W_D,\\
W_D&\rightarrow J_{\mathrm{exp}}\rightarrow A\rightarrow\text{new evidence},\\
\{p_{1:N}^{g,r}\}_{r\in\mathcal T_g}&\rightarrow C_T,\\
Z_T&=(S_t,C_T),\\
\mathcal T_T(W_D,Z_T)&\rightarrow P_T(p_{\mathrm{next}},stop),\\
P_T(p_{\mathrm{next}})&\xrightarrow{\ p_\theta(p\mid S,A)^{-1}\ } A_t.
\end{aligned}}
\]

\begin{plainword}
The whole theory in one chain: observe $\to$ state $\to$ transformation $\to$ mechanisms $\to$ relations $\to$ world $\to$ explore
$\to$ new evidence; and, separately, demonstrations $\to$ task code $\to$ (with the world) a desired next change $\to$ invert to an
action. Every arrow is a claim that must be tested, in order.
\end{plainword}

Its empirical burden is equally concise: show that each higher level contributes stable held-out explanatory or control value that
its lower inputs alone do not, under data with the contrasts needed to make that contribution identifiable.

\begin{example}
The one-line summary of the burden of proof: for every level, show it predicts or controls \emph{new} situations better than the level
below it alone. If the mechanism layer does not beat the transformation layer alone on new data, it has not earned its place.
\end{example}

\section*{References}

\noindent\hangindent=2.5em\hangafter=1 [1] Locatello, F., Weissenborn, D., Unterthiner, T., Mahendran, A., Heigold, G., Uszkoreit, J., Dosovitskiy, A., \& Kipf, T. \textbf{Object-Centric Learning with Slot Attention.} NeurIPS, 2020.\par

\noindent\hangindent=2.5em\hangafter=1 [2] Kipf, T., Elsayed, G. F., Mahendran, A., Stone, A., Sabour, S., Heigold, G., Jonschkowski, R., Dosovitskiy, A., \& Greff, K. \textbf{Conditional Object-Centric Learning from Video.} ICLR, 2022.\par

\noindent\hangindent=2.5em\hangafter=1 [3] Mosbach, M., Ewertz, J. N., Villar-Corrales, A., \& Behnke, S. \textbf{SOLD: Slot Object-Centric Latent Dynamics Models for Relational Manipulation Learning from Pixels.} ICML, 2025.\par

\noindent\hangindent=2.5em\hangafter=1 [4] Dosovitskiy, A., Beyer, L., Kolesnikov, A., Weissenborn, D., Zhai, X., Unterthiner, T., Dehghani, M., Minderer, M., Heigold, G., Gelly, S., Uszkoreit, J., \& Houlsby, N. \textbf{An Image Is Worth 16x16 Words: Transformers for Image Recognition at Scale.} ICLR, 2021.\par

\noindent\hangindent=2.5em\hangafter=1 [5] Sutton, R. S., Precup, D., \& Singh, S. \textbf{Between MDPs and Semi-MDPs: A Framework for Temporal Abstraction in Reinforcement Learning.} Artificial Intelligence, 112(1--2):181--211, 1999.\par

\noindent\hangindent=2.5em\hangafter=1 [6] Shankar, T., \& Gupta, A. \textbf{Learning Robot Skills with Temporal Variational Inference.} ICML, 2020.\par

\noindent\hangindent=2.5em\hangafter=1 [7] Ghaemi, H., Muller, E. B., \& Bakhtiari, S. \textbf{seq-JEPA: Autoregressive Predictive Learning of Invariant-Equivariant World Models.} NeurIPS, 2025.\par

\noindent\hangindent=2.5em\hangafter=1 [8] Nikulin, A., Zisman, I., Tarasov, D., Lyubaykin, N., Polubarov, A., Kiselev, I., \& Kurenkov, V. \textbf{Latent Action Learning Requires Supervision in the Presence of Distractors.} ICML, 2025.\par

\noindent\hangindent=2.5em\hangafter=1 [9] Chen, G., Shao, Q., Cui, T., Zhou, Z., Mao, W., Yang, L., Wang, M., Yang, Y., Chen, H., \& Yue, Y. \textbf{Learning a Unified Latent Action Space from Videos with Action-centric Cycle Consistency.} CVPR, 2026.\par

\noindent\hangindent=2.5em\hangafter=1 [10] Lippe, P., Magliacane, S., L\"owe, S., Asano, Y. M., Cohen, T., \& Gavves, E. \textbf{CITRIS: Causal Identifiability from Temporal Intervened Sequences.} ICML, 2022.\par

\noindent\hangindent=2.5em\hangafter=1 [11] Lippe, P., Magliacane, S., L\"owe, S., Asano, Y. M., Cohen, T., \& Gavves, E. \textbf{BISCUIT: Causal Representation Learning from Binary Interactions.} UAI, 2023.\par

\noindent\hangindent=2.5em\hangafter=1 [12] Lachapelle, S., Rodriguez, P., Sharma, Y., Everett, K. E., Le Priol, R., Lacoste, A., \& Lacoste-Julien, S. \textbf{Disentanglement via Mechanism Sparsity Regularization: A New Principle for Nonlinear ICA.} Conference on Causal Learning and Reasoning, 2022.\par

\noindent\hangindent=2.5em\hangafter=1 [13] Bing, S., Ninad, U., Wahl, J., \& Runge, J. \textbf{Identifying Linearly-Mixed Causal Representations from Multi-Node Interventions.} Conference on Causal Learning and Reasoning, 2024.\par

\noindent\hangindent=2.5em\hangafter=1 [14] Rajendran, G., Reizinger, P., Brendel, W., \& Ravikumar, P. K. \textbf{An Interventional Perspective on Identifiability in Gaussian LTI Systems with Independent Component Analysis.} Conference on Causal Learning and Reasoning, 2024.\par

\noindent\hangindent=2.5em\hangafter=1 [15] Ahuja, K., Mahajan, D., Wang, Y., \& Bengio, Y. \textbf{Interventional Causal Representation Learning.} ICML, 2023.\par

\noindent\hangindent=2.5em\hangafter=1 [16] Lippe, P., Magliacane, S., L\"owe, S., Asano, Y. M., Cohen, T., \& Gavves, E. \textbf{Causal Representation Learning for Instantaneous and Temporal Effects in Interactive Systems.} ICLR, 2023.\par

\noindent\hangindent=2.5em\hangafter=1 [17] Griffiths, T. L., \& Ghahramani, Z. \textbf{Infinite Latent Feature Models and the Indian Buffet Process.} NeurIPS, 2005.\par

\noindent\hangindent=2.5em\hangafter=1 [18] Lee, J., Lee, Y., Kim, J., Kosiorek, A., Choi, S., \& Teh, Y. W. \textbf{Set Transformer: A Framework for Attention-based Permutation-Invariant Neural Networks.} ICML, 2019.\par

\noindent\hangindent=2.5em\hangafter=1 [19] Zaheer, M., Kottur, S., Ravanbakhsh, S., P\'oczos, B., Salakhutdinov, R., \& Smola, A. J. \textbf{Deep Sets.} NeurIPS, 2017.\par

\noindent\hangindent=2.5em\hangafter=1 [20] Louizos, C., Welling, M., \& Kingma, D. P. \textbf{Learning Sparse Neural Networks through $L_0$ Regularization.} ICLR, 2018.\par

\noindent\hangindent=2.5em\hangafter=1 [21] Perez, E., Strub, F., de Vries, H., Dumoulin, V., \& Courville, A. \textbf{FiLM: Visual Reasoning with a General Conditioning Layer.} AAAI, 2018.\par

\noindent\hangindent=2.5em\hangafter=1 [22] Battaglia, P. W., Hamrick, J. B., Bapst, V., et al. \textbf{Relational Inductive Biases, Deep Learning, and Graph Networks.} arXiv:1806.01261, 2018.\par

\noindent\hangindent=2.5em\hangafter=1 [23] Lindley, D. V. \textbf{On a Measure of the Information Provided by an Experiment.} The Annals of Mathematical Statistics, 27(4):986--1005, 1956.\par

\noindent\hangindent=2.5em\hangafter=1 [24] Krause, A., Singh, A., \& Guestrin, C. \textbf{Near-Optimal Sensor Placements in Gaussian Processes: Theory, Efficient Algorithms and Empirical Studies.} Journal of Machine Learning Research, 9:235--284, 2008.\par

\noindent\hangindent=2.5em\hangafter=1 [25] Pathak, D., Agrawal, P., Efros, A. A., \& Darrell, T. \textbf{Curiosity-driven Exploration by Self-supervised Prediction.} ICML, 2017.\par

\noindent\hangindent=2.5em\hangafter=1 [26] Amos, B., Jimenez Rodriguez, I. D., Sacks, J., Boots, B., \& Kolter, J. Z. \textbf{Differentiable MPC for End-to-end Planning and Control.} NeurIPS, 2018.\par

\noindent\hangindent=2.5em\hangafter=1 [27] Chua, K., Calandra, R., McAllister, R., \& Levine, S. \textbf{Deep Reinforcement Learning in a Handful of Trials using Probabilistic Dynamics Models.} NeurIPS, 2018.\par

\noindent\hangindent=2.5em\hangafter=1 [28] Ebert, F., Finn, C., Dasari, S., Xie, A., Lee, A. X., \& Levine, S. \textbf{Visual Foresight: Model-Based Deep Reinforcement Learning for Vision-Based Robotic Control.} arXiv:1812.00568, 2018.\par

\noindent\hangindent=2.5em\hangafter=1 [29] Garrido, Q., Najman, L., \& LeCun, Y. \textbf{Self-supervised Learning of Split Invariant Equivariant Representations.} ICML, 2023.\par

\noindent\hangindent=2.5em\hangafter=1 [30] Freed, B., James, S., Sartoretti, G., \& Choset, H. \textbf{Learning Temporally Abstract World Models without Online Experimentation.} ICML, 2023.\par

\noindent\hangindent=2.5em\hangafter=1 [31] Mishra, N., Abbeel, P., \& Mordatch, I. \textbf{Prediction and Control with Temporal Segment Models.} ICML, 2017.\par

\noindent\hangindent=2.5em\hangafter=1 [32] Locatello, F., Bauer, S., Lucic, M., Raetsch, G., Gelly, S., Sch\"olkopf, B., \& Bachem, O. \textbf{Challenging Common Assumptions in the Unsupervised Learning of Disentangled Representations.} ICML, 2019.\par

\noindent\hangindent=2.5em\hangafter=1 [33] Khemakhem, I., Kingma, D. P., Monti, R. P., \& Hyv\"arinen, A. \textbf{Variational Autoencoders and Nonlinear ICA: A Unifying Framework.} AISTATS, 2020.\par

\noindent\hangindent=2.5em\hangafter=1 [34] Finn, C., Yu, T., Zhang, T., Abbeel, P., \& Levine, S. \textbf{One-Shot Visual Imitation Learning via Meta-Learning.} Conference on Robot Learning, 2017.\par

\noindent\hangindent=2.5em\hangafter=1 [35] James, S., Bloesch, M., \& Davison, A. J. \textbf{Task-Embedded Control Networks for Few-Shot Imitation Learning.} Conference on Robot Learning, 2018.\par

\noindent\hangindent=2.5em\hangafter=1 [36] Rakelly, K., Zhou, A., Finn, C., Levine, S., \& Quillen, D. \textbf{Efficient Off-Policy Meta-Reinforcement Learning via Probabilistic Context Variables.} ICML, 2019.\par

\noindent\hangindent=2.5em\hangafter=1 [37] Garnelo, M., Rosenbaum, D., Maddison, C. J., Ramalho, T., Saxton, D., Shanahan, M., Teh, Y. W., Rezende, D. J., \& Eslami, S. M. A. \textbf{Conditional Neural Processes.} ICML, 2018.\par

\noindent\hangindent=2.5em\hangafter=1 [38] Bishop, C. M. \textbf{Mixture Density Networks.} Aston University Technical Report NCRG/94/004, 1994.\par

\noindent\hangindent=2.5em\hangafter=1 [39] Chi, C., Feng, S., Du, Y., Xu, Z., Cousineau, E., Burchfiel, B., \& Song, S. \textbf{Diffusion Policy: Visuomotor Policy Learning via Action Diffusion.} Robotics: Science and Systems, 2023.\par

\noindent\hangindent=2.5em\hangafter=1 [40] Srinivas, A., Jabri, A., Abbeel, P., Levine, S., \& Finn, C. \textbf{Universal Planning Networks: Learning Generalizable Representations for Visuomotor Control.} ICML, 2018.\par

\end{document}
